\section{Conformance Suite}
\label{sec:conformance}
% ================================================================================

The reference test suite comprises the pytest cases in
\texttt{conformance\_suite.py} (Python 3.10, pytest 9.1.1, Hypothesis
6.168.5; \texttt{requirements.txt}): deterministic tests T1--T158 (some
parametrised over the active cells), and property-based tests
T8--T10 that run random action sequences of up to 40 operations (30 for T10), % count-ok
including a second device, restarts and clock jumps. Every invariant
I1--I15 is exercised (Table~\ref{tab:invariant-independence}). The tests
run the reference \texttt{Bridge} class through its public API only, and
no deterministic test passes by swallowing a refusal
(\texttt{BridgeError}); the property-based driver treats refusals as
expected outcomes and checks the invariants after every action. I8 is
tested through a transport-reported replay flag. I-SIGN-EVENT and
I-SIGN-JWKS concern keys and external services outside the reference
bridge and have no test or vector; I-SIGN-JWKS is a MUST, so a claim
of conformance to it rests on inspection of how the bridge obtains
keys, not on a test, and it has no mutant. I-CHECKPOINT local emission
and field names are tested by T63--T66 and by vectors V52--V54, V88 and
V89, which expect the \texttt{CHECKPOINT} event and run on the Python bridge
only, because the Rust crate has no emission (its runner skips vectors that
require it); external anchoring (SCITT submission) has no test.

\textbf{Important caveat.} The pytest suite tests the
reference \texttt{Bridge} class only and is not an interop test. The
language-neutral vectors (\S\ref{sec:test-vectors}) are the interface for
other implementations; so far they have been run against the Python reference bridge
and the Rust crate, both by the same author; no third-party implementation exists.
Table~\ref{tab:tests} maps the tests to the invariants they exercise;
Table~\ref{tab:traceability} (Appendix~\ref{app:traceability}) traces each
invariant to its TLA+ property, tests, vectors and mutants.

% GENERATED by supplementary/tests_table.py --emit-tex; do not edit.
\begin{longtable}{p{0.12\columnwidth} p{0.2\columnwidth} p{0.58\columnwidth}}
\caption{Tests and the invariants they exercise.}\label{tab:tests}\\
\toprule
\textbf{Tests} & \textbf{Invariants} & \textbf{What they check} \\
\midrule
T1, T7, T10, T11, T19c--T22, T52 & I1, I4, I5, I7, I11 &
Detach emits \texttt{RESUME\_PREPARED} before the cell changes; an
authentic re-attach goes through RESUMING, emits \texttt{RESUME\_VERIFIED}
only on RESUMING $\to$ ATTACHED, and replays exactly the events after the
acknowledged prefix; a history that is not a prefix of the chain is
refused without changing state; a request while RESUMING is refused;
disconnect during resume returns to DETACHED; a broken chain during
resume quarantines the session (T22--T22d); a client resumes however long it was away (T10, T31),
which a held token of limited lifetime could not guarantee. \\
T2, T25 & I2, I13 &
Calls queue under $\textsc{Pend}$ in every attachment, are forwarded only
in $(\textsc{Att},\textsc{None})$, and drain with one
\texttt{TOOL\_CALL\_DRAINED} each. \\
T3, T12, T15, T15b, T56 & I3, I10 &
Widening is rejected in every active cell and emits
\texttt{AUTHORITY\_EXTENSION\_REJECTED}; a widening the runtime announces
itself is reported with source \texttt{observed}, the bridge's record is
unchanged and a revert is requested; narrowing is allowed everywhere as
one event, and a narrowing made while the client was away does not fail
its re-attach. \\
T4, T40 & I6 &
Every mutating operation fails after termination, including across a
restart. \\
T5 & I5 &
A tampered or forged event is detected by chain verification. \\
T6, T26--T29, T30b--T30e & L1, L2, I6 &
Detach and human timeouts, their anchors, simultaneous expiry, and the
zero-timeout edge case; a human wait that outlasts $T_{\text{detach}}$
only while the client is attached; the cumulative cap across episodes,
across restarts and not counting time attached or resuming. \\
T8, T9 & A2, I1--I3, I5, I6, I10, I11, I14, I15 &
Property-based: gap-free sequences, and the invariants hold after every
action of a random sequence, with the cell and the client set equal to
the fold of the log. \\
T13, T13b, T30, T31 & I9, I11 &
An expired nonce is rejected without terminating and without revealing
more than a bad signature; the timer bounds are enforced and derived. \\
T22--T22d, T38 & I5, I12 &
A chain failure after an accepted re-attach, and a recovery failure,
quarantine the session: the failure is recorded, nothing terminates, no
operation or timer acts on it, a restart does not lift it, and a
presenter without a registered key cannot cause it. \\
T57--T59b (vectors V33--V37) & I2, I14 &
A call outside an escalation is refused by default; under the queue
variant it is held until answered after re-attach; under the forward
variant only allowlisted, in-scope permissions are forwarded; while an
escalation is pending nothing is forwarded under any variant. \\
T60--T60d (vectors V38--V40) & I12, I13 &
A drained call awaits delivery and completion; an unacknowledged call is
forwarded again after a restart, with the same key and without a second
drain, only in $(\textsc{Att},\textsc{None})$. \\
T61 & A6, I2 &
Escalations coalesce: a second one is audited, changes no cell and does
not restart $T_{\text{human}}$; one resolution ends them all. \\
T62 (vector V43) & A6 &
A host clock that went backwards across a restart counts as no time
elapsed, and committed timestamps never decrease. \\
T43 (vector V41) & I13, I14 &
An approval given while the client is away lapses after
$T_{\text{approve}}$: the call expires on re-attach and is not
forwarded. \\
T19 (vector V42) & I11 &
A request refused as \texttt{resume\_in\_progress} spends its nonce. \\
T14 & I10 &
\texttt{scope\_version} increments on each strict narrowing; no-op
narrowing is rejected. \\
T16 & I8 &
A transport-reported 0-RTT replay is rejected, cell unchanged. \\
T17--T19b, T51 & I11 &
The signature covers every field; the session, a nonce and the chain (a
bearer) are not enough; a consumed, unknown or foreign-session nonce is
never accepted, and an authentic request spends its nonce even when
rejected later; issuing a challenge writes nothing to the log. \\
T23, T24 & I2, I13 &
Human resolution while detached keeps the queue and drains only after
verification. \\
T32--T40 & I7, I12 &
A failed append changes nothing; nothing is forwarded before its drain
event is durable; restart recovers every cell as detached with an empty
client set, voids outstanding nonces, restarts timers, applies elapsed
timeouts, and fails on a broken chain or a lost key registry. \\
T41--T45 & I13 &
Idempotency keys, queue limit, call expiry, approval expiry, and
revalidation of the current scope at drain; the resume record carries the
queue. \\
T47--T50 & I15 &
A second device joins an attached session; only the last departure
detaches; an added key resumes, an unregistered one is refused, and
registration needs a connected client; a scope narrowed from one device
reaches the other through the replay. \\
T119--T123 (vectors V105--V108) & I15 &
A client registration needs a fresh signature by a connected, unrevoked
key over the operation, session, new key and nonce: unsigned, forged,
other-operation, replayed, expired and revoked-key requests are refused with
a named reason and commit nothing; the signed bytes are pinned by known
answers. \\
T124--T125 & I11 &
The attach proof of possession is the bytes of the tagged tuple that names
the bridge, the session and the key with its scheme tag, pinned by known
answers that a changed field does not verify against; the wire constructors
for attach and client registration carry the scheme. \\
T126 (vector V109) & I11 (a) &
\texttt{challenge} refuses an unknown, a terminated and a quarantined
session with one text that names neither the session nor its state. \\
T127 (vector V110) & I11 (a) &
\texttt{reattach} and \texttt{join} refuse an unknown, a terminated and a
quarantined session with one text that names neither the session nor its
state. \\
T128 (vector V111) & I11-REVOKE &
A key revoked for compromise while its re-attach is verified is not the
attached client afterwards: the re-attach is abandoned (detached, resume
record rewritten) and \texttt{complete\_resume} refuses; a soft revocation
leaves the verification alone. \\
T138--T140 (vectors V121--V126; \texttt{vectors/bytes\_kat.json}) & I-OPERATOR &
An operator signature is bound to the bridge identity and to the log
position (sequence and hash of the last event) at signing: one made for
another bridge, or at an earlier or other position, is refused for a
revocation, a key addition and a termination, and the signed bytes and
signatures are pinned by known answers that a changed field does not verify
against. \\
T141--T142 (vectors V127, V128) & I11, I15 &
A re-attach on an attached session is refused as
\texttt{already\_attached} (rejected, nonce spent, no change of cell) and join
is the way in; the reason is given only to an authentic request, so an
unauthenticated one still gets its own reason. \\
T129--T130 (vectors V112, V113) & I14 &
The approval hash covers the tool and the permission as well as the
content, so approving the arguments of one call does not release another
under a different tool name; the bridge retains the hashed request in the
log and forwards exactly that request, also when it forwards it again after
a restart. \\
T131 (vector V114) & rule Q &
A log that cannot be replayed (no \texttt{ATTACHED} first, or one
without its fields) quarantines its own session only; the other sessions are
recovered and nothing is appended to the malformed log. \\
T132 & I11-REVOKE &
Unsigned revocation requests are written up to a bound per session and
refused past it without a log entry; a signed revocation restores the
budget. \\
T133 & Appendix B &
The reference wire server calls an authentication hook for each unsigned
op and refuses it as \texttt{unauthenticated} on a false answer; signed ops
do not reach the hook. \\
T53--T55 & I14 &
A matching answer drains and the event repeats the hash; a changed view,
a denial or no answer block the forward; the first answer is final. \\
T63--T66, T63b (vectors V52--V54, V88, V89; Python only) & I-CHECKPOINT &
A \texttt{CHECKPOINT} is emitted after the configured number of events
since the last checkpoint and after $T_{\text{cp}}$; its root matches the
recomputed Merkle root; the receipt is optional. The vectors expect the
\texttt{CHECKPOINT} event in the log. \\
T97--T100 (vectors V86, V87) & I11 (f), (g) &
Keyless challenges cannot evict a live nonce and are refused once the
limit is reached; spent and expired nonces make room; unauthenticated
rejections are recorded only up to the bound and are refused uniformly;
the genuine re-attach is accepted after both floods. \\
T101--T104 (vectors V90--V94; \texttt{vectors/bytes\_kat.json}) & I11, I5 &
RFC~8785 known answers (non-ASCII strings, escapes, UTF-16 key order,
numbers, inputs with no canonical form); pinned event and content
hashes; the exact bytes and Ed25519 signature of a re-attach and a join
proof; a proof made for the other operation, for another key id, over
the untagged tuple or without the key id is refused as
\texttt{bad\_signature}. \\
T105--T112, T143 (vectors V95, V96, V102--V104, V129) & I11, I12 &
A key identifier is scoped to its session: a second, validly granted
session that reuses it neither replaces the first session's key nor can
re-attach to the first session, and the owner is not locked out;
\texttt{ATTACHED}, \texttt{CLIENT\_ADDED} and \texttt{OPERATOR\_ADD\_KEY}
carry the scheme tag and public key; recovery repairs the registry from
the log, which is the source of truth: an entry that differs is replaced, one
the log does not record (a crash between the registry write and the event
commit) is dropped, a missing one is written back, and
\texttt{BRIDGE\_RECOVERED} names the keys repaired; a refused append leaves
no registry entry. \\
T113--T118 (vectors V97--V101) & I-GRANT, I11 &
Attach needs a proof of possession by the registered key, under its own
domain tag and bound to the bridge, the session and the key; a missing
proof, a proof by another key and a proof made for another bridge are
refused as \texttt{possession\_invalid}; the grant binds the public key
and the bridge, so a grant seen in transit cannot be used with another
public key under the same key identifier or at another bridge
(\texttt{grant\_invalid}). \\
T67--T69, T75, T76, T86, T87 & I11-REVOKE, I11 &
A soft revocation does not detach an active client; a hard revocation
detaches the revoked key's client and keeps the others connected; a
revoked key is refused on re-attach, also when revoked after its
challenge; two-phase rotation; revoked keys survive a restart. \\
T70--T74, T88, T89 & I11-REVOKE, I6 &
Revoking the last key enters \texttt{key\_exhausted}; client operations are
refused there; the operator adds a recovery key or terminates; the state
and the recovery key survive a restart; refused revocations commit
nothing. \\
T134--T137 (vectors V115--V120) & I11-REVOKE, L1, L2 &
The timers of a \texttt{key\_exhausted} session keep running: it ends
after $T_{\text{detach}}$ (a client operation on an expired session reports
the end, not the exhaustion), an operator key must arrive before the bound,
a restart treats the session like any other and restarts $T_{\text{detach}}$,
a pending escalation times out, and a session whose revoked key is still
connected detaches on transport loss and then ends. \\
T77--T81 & I-GRANT &
An attach is accepted only with a valid grant, which is bound to the
session, scope lineage and key and carries its own domain tag; without
a grant verifier every attach is refused. \\
T82--T85 & I11 &
P-256 proof of possession is accepted and forged proofs are rejected; a
scheme tag is not interchangeable; a known-answer signature from an
independent signer verifies; a server takes only a client's public key. \\
T144 & I11 &
A test-only further scheme tag is accepted at attach through the
pluggable verifier, with no change to the bridge, and refused where the
verifier lacks the tag (extension, Section~\ref{sec:ext-schemes}). \\
T145, T146 (vector V130) & I11-REVOKE, I6 &
An unsigned terminate of a key-exhausted session is refused and commits
nothing; only the operator's signed termination ends it. The refusal
defers to the timers (an expired session is found ended) and to
quarantine (the operator's way out of quarantine stays open). \\
T147 (vector V131) & I11, I15 &
A re-attach naming an unregistered key and one with a bad signature get
the same refusal, so the answer does not list the registered key ids;
the log keeps the true reason. \\
T148--T151 (vector V132) & wire, I-OPERATOR &
The wire server creates a session from an \texttt{attach} message
(refusing a missing grant or proof by name), answers \texttt{resolve} and
\texttt{extend} (the scope in force, a requested revert), verifies the
operator events from the message and names their refusals, and names the
key-exhausted refusal of a bare terminate. \\
T152--T155 (vectors V133, V134) & I11, I11g, I15, wire &
Two sessions on one bridge that hold the same client key share
nothing but the bridge: a nonce issued for one session is unknown to the
other for a re-attach, a join and a registration, and spent by neither
refusal; the budgets of challenges, unauthenticated rejections and
revocation requests are kept per session and are not restored by a
request in another; a tool call key means a call only in its session,
across a restart; every unsigned wire op is asked of the authentication
hook with the message that names its session, and a hook that accepts one
session's caller does not open another. \\
T156--T158 & I11, I14 &
Cases the mutation check could not see while its scratch copies failed
on their own: recovery leaves the registry of a quarantined session as
found; ATTACHED records the scheme of the first key, so a lost entry is
written back with it; a call held while no client is attached carries the
hash of the call a client is shown. \\
T90--T96 (vectors V75--V85) & I-OPERATOR &
Revocation, key addition and termination require a valid operator
signature bound to the session, key and operation; a bridge without an
operator verifier refuses them; the signature is logged and survives
recovery; the operator domain tags separate the signatures. \\
T46 & --- &
The derived view of named states. \\
\bottomrule
\end{longtable}


\paragraph{Mutation check.} A passing suite shows little unless it can
fail. \texttt{mutation\_check.py} applies one-line mutants to the
reference bridge, at least one per invariant I1--I15 and one for the
timer bounds (for example: mutate before append, drop the previous hash,
skip the prefix check, spend a nonce only on success, let an expired
nonce terminate, forward on a mismatching answer hash, detach on a
partial departure, record an observed widening), and requires the suite
to fail on each. A mutant is killed only when pytest exits with code 1
and reports a failed test; a collection, import or syntax error is
reported as an error, not a kill, and a patch that changes nothing must
survive (a control). The script prints the failing tests and whether one
of them carries the mutant's invariant tag in its docstring (kills by a
test with no such tag are listed). Every mutant is killed; the script
exits non-zero otherwise. I-SIGN-JWKS, which is checked by inspection, has none. The mutants include the eviction of a live nonce by a further challenge
and the unbounded logging of unauthenticated failures (I11 (f), (g)), the
dropping of the operation tag or of the key id from the re-attach and join
proof, and, in the canonical encoding, the escaping of non-ASCII text,
code point instead of UTF-16 key order, the printing of an integral float
as \texttt{1.0} and the acceptance of an integer beyond $2^{53}$; and
the count trigger of I-CHECKPOINT
(counting the whole log instead of the events since the last
checkpoint). The Rust crate has a separate, smaller harness,
\texttt{mutation\_check\_rust.py}, which applies one-line mutants to
signature verification, nonce spend, expiry and replay, revoked-key
rejection, the registered-key check, the acknowledged-prefix check, the
hash chain, the operation tag and key id of the resume proof and the
canonical encoding, and requires \texttt{cargo test} to fail on each. It is
a sample of the highest-value rules, not one mutant per invariant; the
mutation evidence for I1--I15 as a whole is on the Python bridge.

\paragraph{Reproducibility.} To reproduce, from
\texttt{paper/supplementary}:

\begin{lstlisting}[language=,basicstyle=\ttfamily\small,frame=single,backgroundcolor=\color{codebg},breaklines=true]
$ python3 -I -m pytest conformance_suite.py -q     # all tests passed
$ python3 -I mutation_check.py                     # every mutant killed
$ python3 -I mutation_check_rust.py                # every Rust mutant killed
$ python3 -I run_vectors.py                        # all vectors passed (V01-V134)
$ python3 -I run_vectors.py --wire                 # wire roundtrip smoke test
$ (cd rust-bridge && cargo test)                   # all vectors on MemStorage and FileStorage
\end{lstlisting}

\noindent The same commands are listed in \texttt{paper/supplementary/REPRODUCE.md}
for readers who prefer a shell-readable copy. The TLC runs are described in
Appendix~\ref{app:formal}. The commands assume a virtual environment in
which \texttt{requirements.txt} is installed (\texttt{python3 -I} ignores
packages installed per user). The runtime requirements of the Python
suite are Python $\geq 3.10$, \texttt{pytest} $\geq 8.0$,
\texttt{hypothesis} $\geq 6.0$, \texttt{cryptography} and
\texttt{jsonschema} (the schema check), with the RFC~8785 encoder bundled
as \texttt{jcs.py}; \texttt{requirements.txt} lists them.

\paragraph{Conformance language and unit.} ``Conformance'' in this
specification is defined at the \emph{bridge} level: a bridge
implementation is conformant if it satisfies all MUST invariants
($I_1, \ldots, I_{15}$, I11-REVOKE, I-GRANT, I-OPERATOR and
I-SIGN-JWKS) for every session it
mediates. Conformance is \emph{not} defined at the session,
deployment, or runtime level; a non-conformant bridge may still
mediate conformant sessions, but its emitted event log is not a
trustworthy audit feed. The tests and vectors provided with this
document check the Python reference bridge and the Rust crate, written by one author; they are not a
third-party interop campaign.

\subsection{Test vectors}\label{sec:test-vectors}

The vectors are language-neutral JSON files in
\texttt{supplementary/vectors/} (\texttt{V01.json}--\texttt{V134.json}).
Each gives a session, an optional configuration, a list of actions and
the expected committed event kinds, final cell and scope version; the
file \texttt{vectors/README.md} defines the schema. The table below is
generated from those files by \texttt{run\_vectors.py --emit-tex} and
checked for staleness by \texttt{run\_vectors.py --check-paper}; a
trailing \texttt{!} marks an action the bridge must refuse. The Python
reference bridge passes all of them (\texttt{python3 -I run\_vectors.py}),
and the Rust crate in \texttt{supplementary/rust-bridge/} runs the same
files, except V52--V54, V88 and V89, which expect checkpoint events the
crate does not emit.

Table~\ref{tab:vectors} lists them.

% GENERATED by supplementary/run_vectors.py --emit-tex; do not edit.
\begingroup\scriptsize\setlength{\tabcolsep}{3pt}
\begin{longtable}{@{}p{0.05\columnwidth} >{\raggedright\arraybackslash}p{0.17\columnwidth} >{\raggedright\arraybackslash}p{0.21\columnwidth} >{\raggedright\arraybackslash}p{0.38\columnwidth} >{\raggedright\arraybackslash}p{0.11\columnwidth}@{}}
\caption{Language-neutral test vectors.}\label{tab:vectors}\\
\toprule
\textbf{ID} & \textbf{Scenario} & \textbf{Actions} & \textbf{Expected event kinds} & \textbf{Final cell} \\
\midrule
\endfirsthead
\toprule
\textbf{ID} & \textbf{Scenario} & \textbf{Actions} & \textbf{Expected event kinds} & \textbf{Final cell} \\
\midrule
\endhead
\texttt{V01} & Detach while a human is pending, then authenticated resume & \texttt{escalate; disconnect; reattach} & \texttt{ATTACHED, HUMAN\_\allowbreak{}PENDING, DETACHMENT, RESUME\_\allowbreak{}PREPARED, RESUME\_\allowbreak{}STARTED, RESUME\_\allowbreak{}VERIFIED} & Att/Pend \\
\texttt{V02} & A rejected re-attach (hash-mismatch) does not terminate; a second re-attach succeeds & \texttt{disconnect; reattach(hash-mismatch)!; reattach} & \texttt{ATTACHED, DETACHMENT, RESUME\_\allowbreak{}PREPARED, RESUME\_\allowbreak{}REJECTED, RESUME\_\allowbreak{}STARTED, RESUME\_\allowbreak{}VERIFIED} & Att/None \\
\texttt{V03} & RESUME\_\allowbreak{}VERIFIED marks Res to Att, not entry to Res & \texttt{disconnect; reattach(no-complete); complete\_\allowbreak{}resume} & \texttt{ATTACHED, DETACHMENT, RESUME\_\allowbreak{}PREPARED, RESUME\_\allowbreak{}STARTED, RESUME\_\allowbreak{}VERIFIED} & Att/None \\
\texttt{V04} & Widening rejected in (Det,Pend) & \texttt{disconnect; escalate; extend([read:docs,exec:shell,exec:sudo])!} & \texttt{ATTACHED, DETACHMENT, RESUME\_\allowbreak{}PREPARED, HUMAN\_\allowbreak{}PENDING, AUTHORITY\_\allowbreak{}EXTENSION\_\allowbreak{}REJECTED} & Det/Pend \\
\texttt{V05} & An expired nonce is rejected and does not terminate; a fresh one resumes & \texttt{disconnect; challenge; advance(60); reattach(n1)!; reattach} & \texttt{ATTACHED, DETACHMENT, RESUME\_\allowbreak{}PREPARED, RESUME\_\allowbreak{}REJECTED, RESUME\_\allowbreak{}STARTED, RESUME\_\allowbreak{}VERIFIED} & Att/None \\
\texttt{V06} & Scope versioning; narrowing only & \texttt{narrow([read:docs]); narrow([]); narrow([read:docs])!} & \texttt{ATTACHED, AUTHORITY\_\allowbreak{}NARROWED, AUTHORITY\_\allowbreak{}NARROWED, AUTHORITY\_\allowbreak{}EXTENSION\_\allowbreak{}REJECTED} & Att/None, sv=2 \\
\texttt{V06b} & The scope is recorded as a sorted set: order and repetition do not change the authority hash & \texttt{disconnect} & \texttt{ATTACHED, DETACHMENT, RESUME\_\allowbreak{}PREPARED} & Det/None \\
\texttt{V06c} & A narrowing that is not strict is refused as not\_\allowbreak{}a\_\allowbreak{}strict\_\allowbreak{}narrowing, whether it repeats the scope or widens it & \texttt{narrow([exec:shell,read:docs,read:docs])!; narrow([read:docs,exec:shell,admin])!; narrow([read:docs])} & \texttt{ATTACHED, AUTHORITY\_\allowbreak{}NARROWING\_\allowbreak{}REJECTED, AUTHORITY\_\allowbreak{}EXTENSION\_\allowbreak{}REJECTED, AUTHORITY\_\allowbreak{}NARROWED} & Att/None, sv=1 \\
\texttt{V07} & A nonce that was used once is never accepted again & \texttt{disconnect; reattach(as:n1); disconnect; reattach(n1)!} & \texttt{ATTACHED, DETACHMENT, RESUME\_\allowbreak{}PREPARED, RESUME\_\allowbreak{}STARTED, RESUME\_\allowbreak{}VERIFIED, DETACHMENT, RESUME\_\allowbreak{}PREPARED, RESUME\_\allowbreak{}REJECTED} & Det/None \\
\texttt{V08} & A request without a valid signature is rejected & \texttt{disconnect; reattach(bad-signature)!} & \texttt{ATTACHED, DETACHMENT, RESUME\_\allowbreak{}PREPARED, RESUME\_\allowbreak{}REJECTED} & Det/None \\
\texttt{V09} & A reported 0-RTT replay is rejected & \texttt{disconnect; reattach(replay)!} & \texttt{ATTACHED, DETACHMENT, RESUME\_\allowbreak{}PREPARED, RESUME\_\allowbreak{}REJECTED} & Det/None \\
\texttt{V09b} & A request the transport flags as a replay has spent its nonce: the same request is refused without the flag & \texttt{disconnect; reattach(as:r1,replay)!; reattach(again:r1)!; reattach} & \texttt{ATTACHED, DETACHMENT, RESUME\_\allowbreak{}PREPARED, RESUME\_\allowbreak{}REJECTED, RESUME\_\allowbreak{}REJECTED, RESUME\_\allowbreak{}STARTED, RESUME\_\allowbreak{}VERIFIED} & Att/None \\
\texttt{V10} & Restart recovers a pending session as detached, then it resumes & \texttt{escalate; tool\_\allowbreak{}call(k1); restart; reattach} & \texttt{ATTACHED, HUMAN\_\allowbreak{}PENDING, TOOL\_\allowbreak{}CALL\_\allowbreak{}QUEUED, BRIDGE\_\allowbreak{}RECOVERED, DETACHMENT, RESUME\_\allowbreak{}PREPARED, RESUME\_\allowbreak{}STARTED, RESUME\_\allowbreak{}VERIFIED} & Att/Pend \\
\texttt{V11} & Restart applies an elapsed human timeout & \texttt{escalate; restart(500)} & \texttt{ATTACHED, HUMAN\_\allowbreak{}PENDING, BRIDGE\_\allowbreak{}RECOVERED, DETACHMENT, RESUME\_\allowbreak{}PREPARED, HUMAN\_\allowbreak{}TIMEOUT, TERMINATION} & Terminated \\
\texttt{V12} & Human resolves while detached: queue drains only after resume (Rule D) & \texttt{disconnect; escalate; tool\_\allowbreak{}call(k1); resolve; reattach} & \texttt{ATTACHED, DETACHMENT, RESUME\_\allowbreak{}PREPARED, HUMAN\_\allowbreak{}PENDING, TOOL\_\allowbreak{}CALL\_\allowbreak{}QUEUED, TOOL\_\allowbreak{}CALL\_\allowbreak{}ANSWERED, HUMAN\_\allowbreak{}RESOLVED, RESUME\_\allowbreak{}STARTED, RESUME\_\allowbreak{}VERIFIED, TOOL\_\allowbreak{}CALL\_\allowbreak{}DRAINED} & Att/None \\
\texttt{V12b} & A resolution with no escalation pending is refused & \texttt{resolve!; disconnect; resolve!} & \texttt{ATTACHED, DETACHMENT, RESUME\_\allowbreak{}PREPARED} & Det/None \\
\texttt{V13} & Queue limit and idempotency keys & \texttt{escalate; tool\_\allowbreak{}call(k1); tool\_\allowbreak{}call(k1); tool\_\allowbreak{}call(k2); tool\_\allowbreak{}call(k3)} & \texttt{ATTACHED, HUMAN\_\allowbreak{}PENDING, TOOL\_\allowbreak{}CALL\_\allowbreak{}QUEUED, TOOL\_\allowbreak{}CALL\_\allowbreak{}QUEUED, TOOL\_\allowbreak{}CALL\_\allowbreak{}REJECTED} & Att/Pend \\
\texttt{V13b} & A generated key does not collide with a key the runtime supplied & \texttt{escalate; tool\_\allowbreak{}call(tc-0002); tool\_\allowbreak{}call; tool\_\allowbreak{}call} & \texttt{ATTACHED, HUMAN\_\allowbreak{}PENDING, TOOL\_\allowbreak{}CALL\_\allowbreak{}QUEUED, TOOL\_\allowbreak{}CALL\_\allowbreak{}QUEUED, TOOL\_\allowbreak{}CALL\_\allowbreak{}QUEUED} & Att/Pend \\
\texttt{V14} & An approval that lapsed before drain is not honoured & \texttt{escalate; disconnect; tool\_\allowbreak{}call(k1); resolve; advance(200); reattach} & \texttt{ATTACHED, HUMAN\_\allowbreak{}PENDING, DETACHMENT, RESUME\_\allowbreak{}PREPARED, TOOL\_\allowbreak{}CALL\_\allowbreak{}QUEUED, TOOL\_\allowbreak{}CALL\_\allowbreak{}ANSWERED, HUMAN\_\allowbreak{}RESOLVED, RESUME\_\allowbreak{}STARTED, RESUME\_\allowbreak{}VERIFIED, TOOL\_\allowbreak{}CALL\_\allowbreak{}EXPIRED} & Att/None \\
\texttt{V15} & Drain revalidates the current scope & \texttt{escalate; tool\_\allowbreak{}call(k1); tool\_\allowbreak{}call(k2); narrow([read:docs]); resolve} & \texttt{ATTACHED, HUMAN\_\allowbreak{}PENDING, TOOL\_\allowbreak{}CALL\_\allowbreak{}QUEUED, TOOL\_\allowbreak{}CALL\_\allowbreak{}QUEUED, AUTHORITY\_\allowbreak{}NARROWED, TOOL\_\allowbreak{}CALL\_\allowbreak{}ANSWERED, TOOL\_\allowbreak{}CALL\_\allowbreak{}ANSWERED, HUMAN\_\allowbreak{}RESOLVED, TOOL\_\allowbreak{}CALL\_\allowbreak{}REJECTED, TOOL\_\allowbreak{}CALL\_\allowbreak{}DRAINED} & Att/None, sv=1 \\
\texttt{V15b} & A direct call outside the current scope is refused and recorded; one inside it is forwarded & \texttt{narrow([read:docs]); tool\_\allowbreak{}call(k1)!; tool\_\allowbreak{}call(k2); tool\_\allowbreak{}call(k3)!} & \texttt{ATTACHED, AUTHORITY\_\allowbreak{}NARROWED, TOOL\_\allowbreak{}CALL\_\allowbreak{}REJECTED, TOOL\_\allowbreak{}CALL\_\allowbreak{}REJECTED} & Att/None, sv=1 \\
\texttt{V16} & Later events do not restart T\_\allowbreak{}detach & \texttt{disconnect; advance(1000); escalate; advance(800); tick} & \texttt{ATTACHED, DETACHMENT, RESUME\_\allowbreak{}PREPARED, HUMAN\_\allowbreak{}PENDING, DETACH\_\allowbreak{}TIMEOUT, TERMINATION} & Terminated \\
\texttt{V16b} & A re-attach in the tick in which T\_\allowbreak{}detach expires is processed first and succeeds & \texttt{disconnect; advance(600,no-tick); reattach} & \texttt{ATTACHED, DETACHMENT, RESUME\_\allowbreak{}PREPARED, RESUME\_\allowbreak{}STARTED, RESUME\_\allowbreak{}VERIFIED} & Att/None \\
\texttt{V16c} & A re-attach refused in the tick in which T\_\allowbreak{}detach expires: the timer fires after it & \texttt{disconnect; advance(600,no-tick); reattach(bad-signature)!} & \texttt{ATTACHED, DETACHMENT, RESUME\_\allowbreak{}PREPARED, RESUME\_\allowbreak{}REJECTED, DETACH\_\allowbreak{}TIMEOUT, TERMINATION} & Terminated \\
\texttt{V16d} & A re-attach after T\_\allowbreak{}detach expired finds the session over: the uniform refusal, and no rejection is logged & \texttt{disconnect; advance(601,no-tick); reattach!} & \texttt{ATTACHED, DETACHMENT, RESUME\_\allowbreak{}PREPARED, DETACH\_\allowbreak{}TIMEOUT, TERMINATION} & Terminated \\
\texttt{V17} & Simultaneous timeouts emit both events & \texttt{escalate; disconnect; advance(100); tick} & \texttt{ATTACHED, HUMAN\_\allowbreak{}PENDING, DETACHMENT, RESUME\_\allowbreak{}PREPARED, HUMAN\_\allowbreak{}TIMEOUT, DETACH\_\allowbreak{}TIMEOUT, TERMINATION} & Terminated \\
\texttt{V18} & Widening rejected in (Res,None) & \texttt{disconnect; reattach(no-complete); extend([read:docs,exec:shell,exec:sudo])!} & \texttt{ATTACHED, DETACHMENT, RESUME\_\allowbreak{}PREPARED, RESUME\_\allowbreak{}STARTED, AUTHORITY\_\allowbreak{}EXTENSION\_\allowbreak{}REJECTED} & Res/None \\
\texttt{V19} & Recovery on a broken hash chain quarantines the session & \texttt{disconnect; tamper\_\allowbreak{}log; restart} & \texttt{ATTACHED, DETACHMENT, RESUME\_\allowbreak{}PREPARED, RECOVERY\_\allowbreak{}FAILED} & Det/None \\
\texttt{V20} & T\_\allowbreak{}call greater than T\_\allowbreak{}human is a configuration error & \texttt{(construct bridge)} & \texttt{---} & configuration error \\
\texttt{V21} & Resume after a scope change made while the client was away & \texttt{disconnect; narrow([read:docs]); reattach} & \texttt{ATTACHED, DETACHMENT, RESUME\_\allowbreak{}PREPARED, AUTHORITY\_\allowbreak{}NARROWED, RESUME\_\allowbreak{}STARTED, RESUME\_\allowbreak{}VERIFIED} & Att/None, sv=1 \\
\texttt{V21b} & The resume record carries the scope and the scope version in force at the detachment & \texttt{narrow([read:docs]); disconnect} & \texttt{ATTACHED, AUTHORITY\_\allowbreak{}NARROWED, DETACHMENT, RESUME\_\allowbreak{}PREPARED} & Det/None, sv=1 \\
\texttt{V22} & A second registered device joins an attached session & \texttt{add\_\allowbreak{}client(ck-2); join(ck-2)} & \texttt{ATTACHED, CLIENT\_\allowbreak{}ADDED, CLIENT\_\allowbreak{}JOINED} & Att/None \\
\texttt{V23} & Only the last departure detaches; the added device resumes & \texttt{add\_\allowbreak{}client(ck-2); join(ck-2); disconnect(ck-1); advance(100000); disconnect(ck-2); reattach(ck-2)} & \texttt{ATTACHED, CLIENT\_\allowbreak{}ADDED, CLIENT\_\allowbreak{}JOINED, CLIENT\_\allowbreak{}LEFT, DETACHMENT, RESUME\_\allowbreak{}PREPARED, RESUME\_\allowbreak{}STARTED, RESUME\_\allowbreak{}VERIFIED} & Att/None \\
\texttt{V23b} & The disconnect of another key does not abandon a re-attach that is being verified & \texttt{add\_\allowbreak{}client(ck-2); disconnect; reattach(ck-1,no-complete); disconnect(ck-2); complete\_\allowbreak{}resume} & \texttt{ATTACHED, CLIENT\_\allowbreak{}ADDED, DETACHMENT, RESUME\_\allowbreak{}PREPARED, RESUME\_\allowbreak{}STARTED, RESUME\_\allowbreak{}VERIFIED} & Att/None \\
\texttt{V24} & A key that was never registered cannot resume & \texttt{disconnect; reattach(ck-9)!} & \texttt{ATTACHED, DETACHMENT, RESUME\_\allowbreak{}PREPARED, RESUME\_\allowbreak{}REJECTED} & Det/None \\
\texttt{V25} & An answer to different content than was queued is not forwarded & \texttt{escalate; tool\_\allowbreak{}call(k1); answer(k1,other); resolve} & \texttt{ATTACHED, HUMAN\_\allowbreak{}PENDING, TOOL\_\allowbreak{}CALL\_\allowbreak{}QUEUED, TOOL\_\allowbreak{}CALL\_\allowbreak{}ANSWERED, HUMAN\_\allowbreak{}RESOLVED, TOOL\_\allowbreak{}CALL\_\allowbreak{}REJECTED} & Att/None \\
\texttt{V26} & A second answer to the same call is refused & \texttt{escalate; tool\_\allowbreak{}call(k1); answer(k1); answer(k1)!; resolve} & \texttt{ATTACHED, HUMAN\_\allowbreak{}PENDING, TOOL\_\allowbreak{}CALL\_\allowbreak{}QUEUED, TOOL\_\allowbreak{}CALL\_\allowbreak{}ANSWERED, HUMAN\_\allowbreak{}RESOLVED, TOOL\_\allowbreak{}CALL\_\allowbreak{}DRAINED} & Att/None \\
\texttt{V27} & A widening announced by the runtime is reported and a revert requested & \texttt{extend([read:docs,exec:shell,exec:sudo],observed)} & \texttt{ATTACHED, AUTHORITY\_\allowbreak{}EXTENSION\_\allowbreak{}REJECTED} & Att/None \\
\texttt{V27b} & A widening the runtime announces is recorded on a key-exhausted session; one asked for is refused as key\_\allowbreak{}exhausted & \texttt{revoke(ck-1,rotation); extend([read:docs,exec:shell,exec:sudo],observed); extend([read:docs,exec:shell,exec:sudo])!} & \texttt{ATTACHED, KEY\_\allowbreak{}REVOKED, AUTHORITY\_\allowbreak{}EXTENSION\_\allowbreak{}REJECTED} & Att/None \\
\texttt{V28} & Persistent profile: resume after six days away & \texttt{disconnect; advance(518400); reattach} & \texttt{ATTACHED, DETACHMENT, RESUME\_\allowbreak{}PREPARED, RESUME\_\allowbreak{}STARTED, RESUME\_\allowbreak{}VERIFIED} & Att/None \\
\texttt{V29} & Chain failure during a re-attach quarantines the session & \texttt{disconnect; tamper\_\allowbreak{}log; reattach!; advance(100000)} & \texttt{ATTACHED, DETACHMENT, RESUME\_\allowbreak{}PREPARED, RESUME\_\allowbreak{}STARTED, RESUME\_\allowbreak{}FAILED} & Res/None \\
\texttt{V30} & T\_\allowbreak{}human may exceed T\_\allowbreak{}detach: the escalation waits while attached, the detach timer ends it once the client is away & \texttt{escalate; advance(2400); disconnect; advance(599); advance(1)} & \texttt{ATTACHED, HUMAN\_\allowbreak{}PENDING, DETACHMENT, RESUME\_\allowbreak{}PREPARED, DETACH\_\allowbreak{}TIMEOUT, TERMINATION} & Terminated \\
\texttt{V31} & The cumulative detachment cap spans episodes & \texttt{disconnect; advance(500); reattach; disconnect; advance(500)} & \texttt{ATTACHED, DETACHMENT, RESUME\_\allowbreak{}PREPARED, RESUME\_\allowbreak{}STARTED, RESUME\_\allowbreak{}VERIFIED, DETACHMENT, RESUME\_\allowbreak{}PREPARED, DETACH\_\allowbreak{}TIMEOUT, TERMINATION} & Terminated \\
\texttt{V31b} & The cumulative cap and a human timeout in one tick: the events name detach\_\allowbreak{}total\_\allowbreak{}timeout & \texttt{disconnect; advance(500); reattach; escalate; disconnect; advance(500)} & \texttt{ATTACHED, DETACHMENT, RESUME\_\allowbreak{}PREPARED, RESUME\_\allowbreak{}STARTED, RESUME\_\allowbreak{}VERIFIED, HUMAN\_\allowbreak{}PENDING, DETACHMENT, RESUME\_\allowbreak{}PREPARED, HUMAN\_\allowbreak{}TIMEOUT, DETACH\_\allowbreak{}TIMEOUT, TERMINATION} & Terminated \\
\texttt{V32} & A crash loop restarts T\_\allowbreak{}detach but not the cumulative cap & \texttt{disconnect; restart(300); restart(300); restart(300); restart(300)} & \texttt{ATTACHED, DETACHMENT, RESUME\_\allowbreak{}PREPARED, BRIDGE\_\allowbreak{}RECOVERED, DETACHMENT, RESUME\_\allowbreak{}PREPARED, BRIDGE\_\allowbreak{}RECOVERED, DETACHMENT, RESUME\_\allowbreak{}PREPARED, BRIDGE\_\allowbreak{}RECOVERED, DETACHMENT, RESUME\_\allowbreak{}PREPARED, BRIDGE\_\allowbreak{}RECOVERED, DETACHMENT, RESUME\_\allowbreak{}PREPARED, DETACH\_\allowbreak{}TIMEOUT, TERMINATION} & Terminated \\
\texttt{V33} & By default a call that arrives away with no escalation is refused & \texttt{disconnect; tool\_\allowbreak{}call(k1)!} & \texttt{ATTACHED, DETACHMENT, RESUME\_\allowbreak{}PREPARED} & Det/None \\
\texttt{V34} & Queue policy: a call away is held, and forwarded only once re-attached and answered & \texttt{disconnect; tool\_\allowbreak{}call(k1); reattach; answer(k1)} & \texttt{ATTACHED, DETACHMENT, RESUME\_\allowbreak{}PREPARED, TOOL\_\allowbreak{}CALL\_\allowbreak{}QUEUED, RESUME\_\allowbreak{}STARTED, RESUME\_\allowbreak{}VERIFIED, TOOL\_\allowbreak{}CALL\_\allowbreak{}ANSWERED, TOOL\_\allowbreak{}CALL\_\allowbreak{}DRAINED} & Att/None \\
\texttt{V35} & Queue policy: an unanswered held call waits, then expires after T\_\allowbreak{}call & \texttt{disconnect; tool\_\allowbreak{}call(k1); reattach; advance(60)} & \texttt{ATTACHED, DETACHMENT, RESUME\_\allowbreak{}PREPARED, TOOL\_\allowbreak{}CALL\_\allowbreak{}QUEUED, RESUME\_\allowbreak{}STARTED, RESUME\_\allowbreak{}VERIFIED, TOOL\_\allowbreak{}CALL\_\allowbreak{}EXPIRED} & Att/None \\
\texttt{V36} & Forward policy: an allowlisted call is forwarded away; any other is refused & \texttt{disconnect; tool\_\allowbreak{}call(k1); tool\_\allowbreak{}call(k2)!} & \texttt{ATTACHED, DETACHMENT, RESUME\_\allowbreak{}PREPARED, TOOL\_\allowbreak{}CALL\_\allowbreak{}QUEUED, TOOL\_\allowbreak{}CALL\_\allowbreak{}ANSWERED, TOOL\_\allowbreak{}CALL\_\allowbreak{}DRAINED} & Det/None \\
\texttt{V37} & Forward policy: while an escalation is pending a call is queued, not forwarded & \texttt{disconnect; escalate; tool\_\allowbreak{}call(k1)} & \texttt{ATTACHED, DETACHMENT, RESUME\_\allowbreak{}PREPARED, HUMAN\_\allowbreak{}PENDING, TOOL\_\allowbreak{}CALL\_\allowbreak{}QUEUED} & Det/Pend \\
\texttt{V38} & A drained call is delivered and completed at most once & \texttt{escalate; tool\_\allowbreak{}call(k1); resolve; delivered(k1); completed(k1); delivered(k1)!} & \texttt{ATTACHED, HUMAN\_\allowbreak{}PENDING, TOOL\_\allowbreak{}CALL\_\allowbreak{}QUEUED, TOOL\_\allowbreak{}CALL\_\allowbreak{}ANSWERED, HUMAN\_\allowbreak{}RESOLVED, TOOL\_\allowbreak{}CALL\_\allowbreak{}DRAINED, TOOL\_\allowbreak{}CALL\_\allowbreak{}DELIVERED, TOOL\_\allowbreak{}CALL\_\allowbreak{}COMPLETED} & Att/None \\
\texttt{V38b} & Completion is accepted once for a call forwarded directly and refused for a call that was denied & \texttt{tool\_\allowbreak{}call(d1); completed(d1); completed(d1)!; escalate; tool\_\allowbreak{}call(k2); answer(k2); resolve; completed(k2)!} & \texttt{ATTACHED, TOOL\_\allowbreak{}CALL\_\allowbreak{}COMPLETED, HUMAN\_\allowbreak{}PENDING, TOOL\_\allowbreak{}CALL\_\allowbreak{}QUEUED, TOOL\_\allowbreak{}CALL\_\allowbreak{}ANSWERED, HUMAN\_\allowbreak{}RESOLVED, TOOL\_\allowbreak{}CALL\_\allowbreak{}REJECTED} & Att/None \\
\texttt{V39} & A drained call that was not delivered is forwarded again after a restart, without a second drain & \texttt{escalate; tool\_\allowbreak{}call(k1); resolve; restart; reattach} & \texttt{ATTACHED, HUMAN\_\allowbreak{}PENDING, TOOL\_\allowbreak{}CALL\_\allowbreak{}QUEUED, TOOL\_\allowbreak{}CALL\_\allowbreak{}ANSWERED, HUMAN\_\allowbreak{}RESOLVED, TOOL\_\allowbreak{}CALL\_\allowbreak{}DRAINED, BRIDGE\_\allowbreak{}RECOVERED, DETACHMENT, RESUME\_\allowbreak{}PREPARED, RESUME\_\allowbreak{}STARTED, RESUME\_\allowbreak{}VERIFIED} & Att/None \\
\texttt{V40} & A delivered call is not forwarded again after a restart & \texttt{escalate; tool\_\allowbreak{}call(k1); resolve; delivered(k1); restart; reattach} & \texttt{ATTACHED, HUMAN\_\allowbreak{}PENDING, TOOL\_\allowbreak{}CALL\_\allowbreak{}QUEUED, TOOL\_\allowbreak{}CALL\_\allowbreak{}ANSWERED, HUMAN\_\allowbreak{}RESOLVED, TOOL\_\allowbreak{}CALL\_\allowbreak{}DRAINED, TOOL\_\allowbreak{}CALL\_\allowbreak{}DELIVERED, BRIDGE\_\allowbreak{}RECOVERED, DETACHMENT, RESUME\_\allowbreak{}PREPARED, RESUME\_\allowbreak{}STARTED, RESUME\_\allowbreak{}VERIFIED} & Att/None \\
\texttt{V41} & An approval given while the client is away lapses after T\_\allowbreak{}approve if the client returns later & \texttt{disconnect; escalate; tool\_\allowbreak{}call(k1); resolve; advance(901); reattach} & \texttt{ATTACHED, DETACHMENT, RESUME\_\allowbreak{}PREPARED, HUMAN\_\allowbreak{}PENDING, TOOL\_\allowbreak{}CALL\_\allowbreak{}QUEUED, TOOL\_\allowbreak{}CALL\_\allowbreak{}ANSWERED, HUMAN\_\allowbreak{}RESOLVED, RESUME\_\allowbreak{}STARTED, RESUME\_\allowbreak{}VERIFIED, TOOL\_\allowbreak{}CALL\_\allowbreak{}EXPIRED} & Att/None \\
\texttt{V41b} & An approval that lapsed is not revived by the resolution of a later escalation & \texttt{escalate; disconnect; tool\_\allowbreak{}call(k1); resolve; advance(150); escalate; tool\_\allowbreak{}call(k2); resolve; advance(50); reattach} & \texttt{ATTACHED, HUMAN\_\allowbreak{}PENDING, DETACHMENT, RESUME\_\allowbreak{}PREPARED, TOOL\_\allowbreak{}CALL\_\allowbreak{}QUEUED, TOOL\_\allowbreak{}CALL\_\allowbreak{}ANSWERED, HUMAN\_\allowbreak{}RESOLVED, HUMAN\_\allowbreak{}PENDING, TOOL\_\allowbreak{}CALL\_\allowbreak{}QUEUED, TOOL\_\allowbreak{}CALL\_\allowbreak{}ANSWERED, HUMAN\_\allowbreak{}RESOLVED, RESUME\_\allowbreak{}STARTED, RESUME\_\allowbreak{}VERIFIED, TOOL\_\allowbreak{}CALL\_\allowbreak{}EXPIRED, TOOL\_\allowbreak{}CALL\_\allowbreak{}DRAINED} & Att/None \\
\texttt{V41c} & An approval is void at exactly T\_\allowbreak{}approve: the call expires & \texttt{escalate; disconnect; tool\_\allowbreak{}call(k1); resolve; advance(100); reattach} & \texttt{ATTACHED, HUMAN\_\allowbreak{}PENDING, DETACHMENT, RESUME\_\allowbreak{}PREPARED, TOOL\_\allowbreak{}CALL\_\allowbreak{}QUEUED, TOOL\_\allowbreak{}CALL\_\allowbreak{}ANSWERED, HUMAN\_\allowbreak{}RESOLVED, RESUME\_\allowbreak{}STARTED, RESUME\_\allowbreak{}VERIFIED, TOOL\_\allowbreak{}CALL\_\allowbreak{}EXPIRED} & Att/None \\
\texttt{V41d} & An approval is valid one second before T\_\allowbreak{}approve: the call is forwarded & \texttt{escalate; disconnect; tool\_\allowbreak{}call(k1); resolve; advance(99); reattach} & \texttt{ATTACHED, HUMAN\_\allowbreak{}PENDING, DETACHMENT, RESUME\_\allowbreak{}PREPARED, TOOL\_\allowbreak{}CALL\_\allowbreak{}QUEUED, TOOL\_\allowbreak{}CALL\_\allowbreak{}ANSWERED, HUMAN\_\allowbreak{}RESOLVED, RESUME\_\allowbreak{}STARTED, RESUME\_\allowbreak{}VERIFIED, TOOL\_\allowbreak{}CALL\_\allowbreak{}DRAINED} & Att/None \\
\texttt{V42} & A request refused as resume\_\allowbreak{}in\_\allowbreak{}progress spends its nonce & \texttt{disconnect; challenge; reattach(no-complete); reattach(n2)!; complete\_\allowbreak{}resume; disconnect; reattach(n2)!} & \texttt{ATTACHED, DETACHMENT, RESUME\_\allowbreak{}PREPARED, RESUME\_\allowbreak{}STARTED, RESUME\_\allowbreak{}REJECTED, RESUME\_\allowbreak{}VERIFIED, DETACHMENT, RESUME\_\allowbreak{}PREPARED, RESUME\_\allowbreak{}REJECTED} & Det/None \\
\texttt{V43} & A clock that went backwards across a restart counts as no time elapsed & \texttt{disconnect; advance(500); reattach; disconnect; restart(-1000); advance(1999); advance(1)} & \texttt{ATTACHED, DETACHMENT, RESUME\_\allowbreak{}PREPARED, RESUME\_\allowbreak{}STARTED, RESUME\_\allowbreak{}VERIFIED, DETACHMENT, RESUME\_\allowbreak{}PREPARED, BRIDGE\_\allowbreak{}RECOVERED, DETACHMENT, RESUME\_\allowbreak{}PREPARED, DETACH\_\allowbreak{}TIMEOUT, TERMINATION} & Terminated \\
\texttt{V44} & Wire: simplest detach and reattach (no escalation) & \texttt{disconnect; reattach} & \texttt{ATTACHED, DETACHMENT, RESUME\_\allowbreak{}PREPARED, RESUME\_\allowbreak{}STARTED, RESUME\_\allowbreak{}VERIFIED} & Att/None \\
\texttt{V45} & Wire: add\_\allowbreak{}client then join with the new key (multi-device) & \texttt{add\_\allowbreak{}client(ck-2); join(ck-2)} & \texttt{ATTACHED, CLIENT\_\allowbreak{}ADDED, CLIENT\_\allowbreak{}JOINED} & Att/None \\
\texttt{V46} & Wire: tool\_\allowbreak{}call with explicit content, answer, drain & \texttt{escalate; tool\_\allowbreak{}call(k1); answer(k1); resolve} & \texttt{ATTACHED, HUMAN\_\allowbreak{}PENDING, TOOL\_\allowbreak{}CALL\_\allowbreak{}QUEUED, TOOL\_\allowbreak{}CALL\_\allowbreak{}ANSWERED, HUMAN\_\allowbreak{}RESOLVED, TOOL\_\allowbreak{}CALL\_\allowbreak{}DRAINED} & Att/None \\
\texttt{V47} & Wire: narrow then extend (observed) — authority wire ops & \texttt{narrow([read:docs]); extend([read:docs,exec:shell,exec:sudo],observed)} & \texttt{ATTACHED, AUTHORITY\_\allowbreak{}NARROWED, AUTHORITY\_\allowbreak{}EXTENSION\_\allowbreak{}REJECTED} & Att/None, sv=1 \\
\texttt{V48} & Wire: delivered then completed — forwarding acknowledgement chain & \texttt{escalate; tool\_\allowbreak{}call(k1); answer(k1); resolve; delivered(k1); completed(k1)} & \texttt{ATTACHED, HUMAN\_\allowbreak{}PENDING, TOOL\_\allowbreak{}CALL\_\allowbreak{}QUEUED, TOOL\_\allowbreak{}CALL\_\allowbreak{}ANSWERED, HUMAN\_\allowbreak{}RESOLVED, TOOL\_\allowbreak{}CALL\_\allowbreak{}DRAINED, TOOL\_\allowbreak{}CALL\_\allowbreak{}DELIVERED, TOOL\_\allowbreak{}CALL\_\allowbreak{}COMPLETED} & Att/None \\
\texttt{V49} & Wire: terminate from (Att, None) — terminate wire op & \texttt{terminate(clean\_\allowbreak{}end)} & \texttt{ATTACHED, TERMINATION} & Terminated \\
\texttt{V50} & Wire: tick fires detach timeout — tick wire op & \texttt{disconnect; advance(1800); tick} & \texttt{ATTACHED, DETACHMENT, RESUME\_\allowbreak{}PREPARED, DETACH\_\allowbreak{}TIMEOUT, TERMINATION} & Terminated \\
\texttt{V51} & Wire: disconnect specific key, then disconnect all — partial disconnect & \texttt{add\_\allowbreak{}client(ck-2); join(ck-2); disconnect(ck-1); disconnect} & \texttt{ATTACHED, CLIENT\_\allowbreak{}ADDED, CLIENT\_\allowbreak{}JOINED, CLIENT\_\allowbreak{}LEFT, DETACHMENT, RESUME\_\allowbreak{}PREPARED} & Det/None \\
\texttt{V52} & Checkpoint emitted by the event threshold, with a tool call queued & \texttt{escalate; tool\_\allowbreak{}call(k1)} & \texttt{ATTACHED, HUMAN\_\allowbreak{}PENDING, TOOL\_\allowbreak{}CALL\_\allowbreak{}QUEUED, CHECKPOINT} & Att/Pend \\
\texttt{V53} & Checkpoint emitted by the time threshold on a tick & \texttt{escalate; advance(100); tick} & \texttt{ATTACHED, HUMAN\_\allowbreak{}PENDING, CHECKPOINT} & Att/Pend \\
\texttt{V54} & A checkpointed log that is tampered with fails recovery (rule Q) & \texttt{escalate; tool\_\allowbreak{}call(k1); tamper\_\allowbreak{}log; restart} & \texttt{ATTACHED, HUMAN\_\allowbreak{}PENDING, TOOL\_\allowbreak{}CALL\_\allowbreak{}QUEUED, CHECKPOINT, RECOVERY\_\allowbreak{}FAILED} & Att/Pend \\
\texttt{V54b} & No checkpoint follows a TERMINATION & \texttt{escalate; terminate} & \texttt{ATTACHED, HUMAN\_\allowbreak{}PENDING, CHECKPOINT, TERMINATION} & Terminated \\
\texttt{V54c} & No checkpoint over a log that does not verify, nor on the session it quarantines & \texttt{disconnect; tamper\_\allowbreak{}log; reattach!; advance(100)} & \texttt{ATTACHED, DETACHMENT, RESUME\_\allowbreak{}PREPARED, CHECKPOINT, RESUME\_\allowbreak{}STARTED, RESUME\_\allowbreak{}FAILED} & Res/None \\
\texttt{V55} & Attach without a scope grant is refused & \texttt{attach with grant none} & \texttt{(none: attach refused, grant\_\allowbreak{}invalid)} & no session \\
\texttt{V56} & Attach with a grant from another authority is refused & \texttt{attach with grant forged} & \texttt{(none: attach refused, grant\_\allowbreak{}invalid)} & no session \\
\texttt{V57} & Attach with a grant issued for another session is refused & \texttt{attach with grant other-session} & \texttt{(none: attach refused, grant\_\allowbreak{}invalid)} & no session \\
\texttt{V58} & Attach claiming a wider scope than the grant is refused & \texttt{attach with grant wider-scope} & \texttt{(none: attach refused, grant\_\allowbreak{}invalid)} & no session \\
\texttt{V59} & P-256: a valid signature under the registered key resumes & \texttt{disconnect; reattach} & \texttt{ATTACHED, DETACHMENT, RESUME\_\allowbreak{}PREPARED, RESUME\_\allowbreak{}STARTED, RESUME\_\allowbreak{}VERIFIED} & Att/None \\
\texttt{V60} & P-256: a real signature by another key is rejected & \texttt{disconnect; reattach(wrong-key)!} & \texttt{ATTACHED, DETACHMENT, RESUME\_\allowbreak{}PREPARED, RESUME\_\allowbreak{}REJECTED} & Det/None \\
\texttt{V61} & P-256: a valid signature over a different message is rejected & \texttt{disconnect; reattach(other-message)!} & \texttt{ATTACHED, DETACHMENT, RESUME\_\allowbreak{}PREPARED, RESUME\_\allowbreak{}REJECTED} & Det/None \\
\texttt{V62} & P-256: a nonce issued for a key before it was registered never verifies, also after the key is added & \texttt{challenge(ck-2); add\_\allowbreak{}client(ck-2); disconnect; reattach(ck-2,n1)!; reattach(ck-2)} & \texttt{ATTACHED, CLIENT\_\allowbreak{}ADDED, DETACHMENT, RESUME\_\allowbreak{}PREPARED, RESUME\_\allowbreak{}REJECTED, RESUME\_\allowbreak{}STARTED, RESUME\_\allowbreak{}VERIFIED} & Att/None \\
\texttt{V62b} & A join refused before its key was registered is refused again, byte for byte, after the registration & \texttt{join(ck-2,as:r1)!; add\_\allowbreak{}client(ck-2); join(again:r1)!; join(ck-2,r1)!; join(ck-2)} & \texttt{ATTACHED, RESUME\_\allowbreak{}REJECTED, CLIENT\_\allowbreak{}ADDED, RESUME\_\allowbreak{}REJECTED, RESUME\_\allowbreak{}REJECTED, CLIENT\_\allowbreak{}JOINED} & Att/None \\
\texttt{V62c} & A re-attach refused on a key-exhausted session is refused again after the operator adds its key & \texttt{revoke(ck-1,compromise); reattach(ck-2,as:r1)!; reattach(ck-1)!; operator\_\allowbreak{}add\_\allowbreak{}key(ck-2); reattach(again:r1)!; reattach(ck-2)} & \texttt{ATTACHED, KEY\_\allowbreak{}REVOKED, DETACHMENT, RESUME\_\allowbreak{}PREPARED, RESUME\_\allowbreak{}REJECTED, RESUME\_\allowbreak{}REJECTED, OPERATOR\_\allowbreak{}ADD\_\allowbreak{}KEY, RESUME\_\allowbreak{}REJECTED, RESUME\_\allowbreak{}STARTED, RESUME\_\allowbreak{}VERIFIED} & Att/None \\
\texttt{V62d} & A nonce issued for one key is refused for another registered key and stays live for its own & \texttt{add\_\allowbreak{}client(ck-2); disconnect; challenge(ck-1); reattach(ck-2,n1)!; reattach(ck-1,n1)} & \texttt{ATTACHED, CLIENT\_\allowbreak{}ADDED, DETACHMENT, RESUME\_\allowbreak{}PREPARED, RESUME\_\allowbreak{}REJECTED, RESUME\_\allowbreak{}STARTED, RESUME\_\allowbreak{}VERIFIED} & Att/None \\
\texttt{V62e} & A challenge is answered for a key id the session does not register, and counts against the same limit & \texttt{disconnect; challenge(nobody,16); challenge(ck-1)!; challenge(nobody)!} & \texttt{ATTACHED, DETACHMENT, RESUME\_\allowbreak{}PREPARED} & Det/None \\
\texttt{V62f} & A client registration with a nonce issued for another key is refused as a bad signature & \texttt{add\_\allowbreak{}client(ck-2); challenge(ck-2); add\_\allowbreak{}client(ck-3,bad\_\allowbreak{}signature,n1)!; add\_\allowbreak{}client(ck-3)} & \texttt{ATTACHED, CLIENT\_\allowbreak{}ADDED, CLIENT\_\allowbreak{}ADDED} & Att/None \\
\texttt{V63} & Soft revocation keeps the connected client & \texttt{add\_\allowbreak{}client(ck-2); revoke(ck-1,rotation)} & \texttt{ATTACHED, CLIENT\_\allowbreak{}ADDED, KEY\_\allowbreak{}REVOKED} & Att/None \\
\texttt{V64} & Two-phase rotation: add the new key, revoke the old, only the new one re-attaches & \texttt{add\_\allowbreak{}client(ck-2); revoke(ck-1,rotation); disconnect; reattach(ck-1)!; reattach(ck-2)} & \texttt{ATTACHED, CLIENT\_\allowbreak{}ADDED, KEY\_\allowbreak{}REVOKED, DETACHMENT, RESUME\_\allowbreak{}PREPARED, RESUME\_\allowbreak{}REJECTED, RESUME\_\allowbreak{}STARTED, RESUME\_\allowbreak{}VERIFIED} & Att/None \\
\texttt{V65} & Hard revocation removes only the revoked client while another stays connected & \texttt{add\_\allowbreak{}client(ck-2); join(ck-2); revoke(ck-1,compromise)} & \texttt{ATTACHED, CLIENT\_\allowbreak{}ADDED, CLIENT\_\allowbreak{}JOINED, KEY\_\allowbreak{}REVOKED} & Att/None \\
\texttt{V66} & Revoking the last key by compromise detaches and leaves the session key exhausted; its re-attach is refused as revoked & \texttt{revoke(ck-1,compromise); escalate!; reattach(ck-1)!; revoke(ck-1,compromise)!} & \texttt{ATTACHED, KEY\_\allowbreak{}REVOKED, DETACHMENT, RESUME\_\allowbreak{}PREPARED, RESUME\_\allowbreak{}REJECTED} & Det/None \\
\texttt{V67} & OPERATOR\_\allowbreak{}ADD\_\allowbreak{}KEY recovers a key exhausted session; the revoked key stays revoked & \texttt{revoke(ck-1,compromise); operator\_\allowbreak{}add\_\allowbreak{}key(ck-2); reattach(ck-2); disconnect; reattach(ck-1)!} & \texttt{ATTACHED, KEY\_\allowbreak{}REVOKED, DETACHMENT, RESUME\_\allowbreak{}PREPARED, OPERATOR\_\allowbreak{}ADD\_\allowbreak{}KEY, RESUME\_\allowbreak{}STARTED, RESUME\_\allowbreak{}VERIFIED, DETACHMENT, RESUME\_\allowbreak{}PREPARED, RESUME\_\allowbreak{}REJECTED} & Det/None \\
\texttt{V68} & OPERATOR\_\allowbreak{}TERMINATION ends a key exhausted session & \texttt{revoke(ck-1,compromise); operator\_\allowbreak{}terminate; operator\_\allowbreak{}add\_\allowbreak{}key(ck-2)!} & \texttt{ATTACHED, KEY\_\allowbreak{}REVOKED, DETACHMENT, RESUME\_\allowbreak{}PREPARED, OPERATOR\_\allowbreak{}TERMINATION, TERMINATION} & Terminated \\
\texttt{V69} & A client's revocation request is recorded and revokes nothing & \texttt{revoke(ck-1,rotation,requested,unsigned)} & \texttt{ATTACHED, KEY\_\allowbreak{}REVOCATION\_\allowbreak{}REQUESTED} & Att/None \\
\texttt{V70} & A key revoked between its challenge and its re-attach is rejected, and no other key can use that nonce & \texttt{add\_\allowbreak{}client(ck-2); disconnect; challenge; revoke(ck-1,rotation); reattach(ck-1,n1)!; reattach(ck-2,n1)!; reattach(ck-2)} & \texttt{ATTACHED, CLIENT\_\allowbreak{}ADDED, DETACHMENT, RESUME\_\allowbreak{}PREPARED, KEY\_\allowbreak{}REVOKED, RESUME\_\allowbreak{}REJECTED, RESUME\_\allowbreak{}REJECTED, RESUME\_\allowbreak{}STARTED, RESUME\_\allowbreak{}VERIFIED} & Att/None \\
\texttt{V70b} & Refusals of a revoked key are logged only up to the bound, then the active key re-attaches & \texttt{add\_\allowbreak{}client(ck-2); disconnect; revoke(ck-1,rotation); reattach(ck-1,unknown,20)!; reattach(ck-2)} & \texttt{ATTACHED, CLIENT\_\allowbreak{}ADDED, DETACHMENT, RESUME\_\allowbreak{}PREPARED, KEY\_\allowbreak{}REVOKED, RESUME\_\allowbreak{}REJECTED, RESUME\_\allowbreak{}REJECTED, RESUME\_\allowbreak{}REJECTED, RESUME\_\allowbreak{}REJECTED, RESUME\_\allowbreak{}REJECTED, RESUME\_\allowbreak{}REJECTED, RESUME\_\allowbreak{}REJECTED, RESUME\_\allowbreak{}REJECTED, RESUME\_\allowbreak{}REJECTED, RESUME\_\allowbreak{}REJECTED, RESUME\_\allowbreak{}REJECTED, RESUME\_\allowbreak{}REJECTED, RESUME\_\allowbreak{}REJECTED, RESUME\_\allowbreak{}REJECTED, RESUME\_\allowbreak{}REJECTED, RESUME\_\allowbreak{}REJECTED, RESUME\_\allowbreak{}STARTED, RESUME\_\allowbreak{}VERIFIED} & Att/None \\
\texttt{V71} & A revoked key stays revoked across a restart & \texttt{add\_\allowbreak{}client(ck-2); revoke(ck-1,rotation); restart; reattach(ck-1)!; reattach(ck-2)} & \texttt{ATTACHED, CLIENT\_\allowbreak{}ADDED, KEY\_\allowbreak{}REVOKED, BRIDGE\_\allowbreak{}RECOVERED, DETACHMENT, RESUME\_\allowbreak{}PREPARED, RESUME\_\allowbreak{}REJECTED, RESUME\_\allowbreak{}STARTED, RESUME\_\allowbreak{}VERIFIED} & Att/None \\
\texttt{V72} & Key exhaustion and the recovery key are rebuilt from the log after restarts & \texttt{revoke(ck-1,compromise); restart; escalate!; operator\_\allowbreak{}add\_\allowbreak{}key(ck-2); restart; reattach(ck-2)} & \texttt{ATTACHED, KEY\_\allowbreak{}REVOKED, DETACHMENT, RESUME\_\allowbreak{}PREPARED, BRIDGE\_\allowbreak{}RECOVERED, DETACHMENT, RESUME\_\allowbreak{}PREPARED, OPERATOR\_\allowbreak{}ADD\_\allowbreak{}KEY, BRIDGE\_\allowbreak{}RECOVERED, DETACHMENT, RESUME\_\allowbreak{}PREPARED, RESUME\_\allowbreak{}STARTED, RESUME\_\allowbreak{}VERIFIED} & Att/None \\
\texttt{V73} & Revocation refusals: unknown key, add-key outside exhaustion, a registered key id, revoking when exhausted & \texttt{revoke(ck-9)!; operator\_\allowbreak{}add\_\allowbreak{}key(ck-2)!; operator\_\allowbreak{}terminate!; revoke(ck-1,compromise); operator\_\allowbreak{}add\_\allowbreak{}key(ck-1)!; revoke(ck-1,compromise)!} & \texttt{ATTACHED, KEY\_\allowbreak{}REVOKED, DETACHMENT, RESUME\_\allowbreak{}PREPARED} & Det/None \\
\texttt{V73b} & A revocation names rotation or compromise, and a key that is revoked is not revoked again & \texttt{add\_\allowbreak{}client(ck-2); revoke(ck-1,lost)!; revoke(ck-1,rotation); revoke(ck-1,compromise)!; revoke(ck-1,rotation,requested,unsigned)!} & \texttt{ATTACHED, CLIENT\_\allowbreak{}ADDED, KEY\_\allowbreak{}REVOKED} & Att/None \\
\texttt{V74} & Soft revocation of the only connected key keeps it in the client set until it disconnects, then it cannot re-attach & \texttt{add\_\allowbreak{}client(ck-2); revoke(ck-1,rotation); disconnect(ck-1); reattach(ck-1)!} & \texttt{ATTACHED, CLIENT\_\allowbreak{}ADDED, KEY\_\allowbreak{}REVOKED, DETACHMENT, RESUME\_\allowbreak{}PREPARED, RESUME\_\allowbreak{}REJECTED} & Det/None \\
\texttt{V75} & A revocation signed by the operator is accepted and logged & \texttt{revoke(ck-1,rotation)} & \texttt{ATTACHED, KEY\_\allowbreak{}REVOKED} & Att/None \\
\texttt{V76} & A revocation with an operator signature missing is refused and commits nothing & \texttt{revoke(ck-1,compromise,sig-none)!} & \texttt{ATTACHED} & Att/None \\
\texttt{V77} & A revocation with an operator signature signed by another key is refused and commits nothing & \texttt{revoke(ck-1,compromise,sig-forged)!} & \texttt{ATTACHED} & Att/None \\
\texttt{V78} & A revocation with an operator signature signed over another operation is refused and commits nothing & \texttt{revoke(ck-1,compromise,sig-other-operation)!} & \texttt{ATTACHED} & Att/None \\
\texttt{V79} & An operator key addition with a signature missing is refused and the session stays key exhausted & \texttt{revoke(ck-1,compromise); operator\_\allowbreak{}add\_\allowbreak{}key(ck-2,sig-none)!} & \texttt{ATTACHED, KEY\_\allowbreak{}REVOKED, DETACHMENT, RESUME\_\allowbreak{}PREPARED} & Det/None \\
\texttt{V80} & An operator key addition with a signature signed by another key is refused and the session stays key exhausted & \texttt{revoke(ck-1,compromise); operator\_\allowbreak{}add\_\allowbreak{}key(ck-2,sig-forged)!} & \texttt{ATTACHED, KEY\_\allowbreak{}REVOKED, DETACHMENT, RESUME\_\allowbreak{}PREPARED} & Det/None \\
\texttt{V81} & An operator key addition with a signature signed over another operation is refused and the session stays key exhausted & \texttt{revoke(ck-1,compromise); operator\_\allowbreak{}add\_\allowbreak{}key(ck-2,sig-other-operation)!} & \texttt{ATTACHED, KEY\_\allowbreak{}REVOKED, DETACHMENT, RESUME\_\allowbreak{}PREPARED} & Det/None \\
\texttt{V82} & An operator termination with a signature missing is refused and the session stays key exhausted & \texttt{revoke(ck-1,compromise); operator\_\allowbreak{}terminate(sig-none)!} & \texttt{ATTACHED, KEY\_\allowbreak{}REVOKED, DETACHMENT, RESUME\_\allowbreak{}PREPARED} & Det/None \\
\texttt{V83} & An operator termination with a signature signed by another key is refused and the session stays key exhausted & \texttt{revoke(ck-1,compromise); operator\_\allowbreak{}terminate(sig-forged)!} & \texttt{ATTACHED, KEY\_\allowbreak{}REVOKED, DETACHMENT, RESUME\_\allowbreak{}PREPARED} & Det/None \\
\texttt{V84} & An operator termination with a signature signed over another operation is refused and the session stays key exhausted & \texttt{revoke(ck-1,compromise); operator\_\allowbreak{}terminate(sig-other-operation)!} & \texttt{ATTACHED, KEY\_\allowbreak{}REVOKED, DETACHMENT, RESUME\_\allowbreak{}PREPARED} & Det/None \\
\texttt{V85} & A signed operator termination ends a key-exhausted session & \texttt{revoke(ck-1,compromise); operator\_\allowbreak{}terminate} & \texttt{ATTACHED, KEY\_\allowbreak{}REVOKED, DETACHMENT, RESUME\_\allowbreak{}PREPARED, OPERATOR\_\allowbreak{}TERMINATION, TERMINATION} & Terminated \\
\texttt{V86} & Keyless challenges cannot evict the nonce a genuine client holds & \texttt{disconnect; challenge; challenge(15); challenge!; reattach(genuine)} & \texttt{ATTACHED, DETACHMENT, RESUME\_\allowbreak{}PREPARED, RESUME\_\allowbreak{}STARTED, RESUME\_\allowbreak{}VERIFIED} & Att/None \\
\texttt{V87} & Unauthenticated rejections are logged only up to a bound, then the genuine re-attach succeeds & \texttt{disconnect; reattach(unknown,20,bad-signature)!; reattach} & \texttt{ATTACHED, DETACHMENT, RESUME\_\allowbreak{}PREPARED, RESUME\_\allowbreak{}REJECTED, RESUME\_\allowbreak{}REJECTED, RESUME\_\allowbreak{}REJECTED, RESUME\_\allowbreak{}REJECTED, RESUME\_\allowbreak{}REJECTED, RESUME\_\allowbreak{}REJECTED, RESUME\_\allowbreak{}REJECTED, RESUME\_\allowbreak{}REJECTED, RESUME\_\allowbreak{}REJECTED, RESUME\_\allowbreak{}REJECTED, RESUME\_\allowbreak{}REJECTED, RESUME\_\allowbreak{}REJECTED, RESUME\_\allowbreak{}REJECTED, RESUME\_\allowbreak{}REJECTED, RESUME\_\allowbreak{}REJECTED, RESUME\_\allowbreak{}REJECTED, RESUME\_\allowbreak{}STARTED, RESUME\_\allowbreak{}VERIFIED} & Att/None \\
\texttt{V88} & Checkpoint emitted when the event threshold is reached & \texttt{escalate; tool\_\allowbreak{}call(k1)} & \texttt{ATTACHED, HUMAN\_\allowbreak{}PENDING, TOOL\_\allowbreak{}CALL\_\allowbreak{}QUEUED, CHECKPOINT} & Att/Pend \\
\texttt{V89} & Checkpoint emitted when the time threshold is reached & \texttt{escalate; advance(100); tick} & \texttt{ATTACHED, HUMAN\_\allowbreak{}PENDING, CHECKPOINT} & Att/Pend \\
\texttt{V90} & A signature made for a re-attach is refused when presented as a join & \texttt{add\_\allowbreak{}client(ck-2); join(ck-2,other-operation)!} & \texttt{ATTACHED, CLIENT\_\allowbreak{}ADDED, RESUME\_\allowbreak{}REJECTED} & Att/None \\
\texttt{V90b} & A join on a detached session: unverified it is a bad signature, verified it is refused as not\_\allowbreak{}attached and starts no re-attach & \texttt{add\_\allowbreak{}client(ck-2); disconnect; join(ck-2,bad-signature)!; join(ck-2,as:r1)!; join(again:r1)!; reattach(ck-2)} & \texttt{ATTACHED, CLIENT\_\allowbreak{}ADDED, DETACHMENT, RESUME\_\allowbreak{}PREPARED, RESUME\_\allowbreak{}REJECTED, RESUME\_\allowbreak{}REJECTED, RESUME\_\allowbreak{}REJECTED, RESUME\_\allowbreak{}STARTED, RESUME\_\allowbreak{}VERIFIED} & Att/None \\
\texttt{V90c} & A revoked key cannot join & \texttt{add\_\allowbreak{}client(ck-2); revoke(ck-2,rotation); join(ck-2)!} & \texttt{ATTACHED, CLIENT\_\allowbreak{}ADDED, KEY\_\allowbreak{}REVOKED, RESUME\_\allowbreak{}REJECTED} & Att/None \\
\texttt{V90d} & A join while a re-attach is being verified: unverified it is a bad signature, verified it is refused as resume\_\allowbreak{}in\_\allowbreak{}progress & \texttt{add\_\allowbreak{}client(ck-2); disconnect; reattach(ck-1,no-complete); join(ck-2,bad-signature)!; join(ck-2)!; complete\_\allowbreak{}resume; join(ck-2)} & \texttt{ATTACHED, CLIENT\_\allowbreak{}ADDED, DETACHMENT, RESUME\_\allowbreak{}PREPARED, RESUME\_\allowbreak{}STARTED, RESUME\_\allowbreak{}REJECTED, RESUME\_\allowbreak{}REJECTED, RESUME\_\allowbreak{}VERIFIED, CLIENT\_\allowbreak{}JOINED} & Att/None \\
\texttt{V91} & A signature made for a join is refused when presented as a re-attach & \texttt{disconnect; reattach(other-operation)!} & \texttt{ATTACHED, DETACHMENT, RESUME\_\allowbreak{}PREPARED, RESUME\_\allowbreak{}REJECTED} & Det/None \\
\texttt{V92} & A re-attach proof over the untagged tuple of an earlier format is refused & \texttt{disconnect; reattach!} & \texttt{ATTACHED, DETACHMENT, RESUME\_\allowbreak{}PREPARED, RESUME\_\allowbreak{}REJECTED} & Det/None \\
\texttt{V93} & A re-attach proof that omits the signing key id is refused & \texttt{disconnect; reattach!} & \texttt{ATTACHED, DETACHMENT, RESUME\_\allowbreak{}PREPARED, RESUME\_\allowbreak{}REJECTED} & Det/None \\
\texttt{V94} & P-256: a signature made for a join is refused when presented as a re-attach & \texttt{disconnect; reattach(other-operation)!} & \texttt{ATTACHED, DETACHMENT, RESUME\_\allowbreak{}PREPARED, RESUME\_\allowbreak{}REJECTED} & Det/None \\
\texttt{V95} & A second session that reuses the key id leaves the first session's key in place and the owner re-attaches & \texttt{attach\_\allowbreak{}other(ck-1,sess-95b); disconnect; reattach(ck-1)} & \texttt{ATTACHED, DETACHMENT, RESUME\_\allowbreak{}PREPARED, RESUME\_\allowbreak{}STARTED, RESUME\_\allowbreak{}VERIFIED} & Att/None \\
\texttt{V96} & A second session that reuses the key id cannot re-attach to the first session with its own key & \texttt{attach\_\allowbreak{}other(ck-1,sess-96b); disconnect; reattach(ck-1,other-session-key)!} & \texttt{ATTACHED, DETACHMENT, RESUME\_\allowbreak{}PREPARED, RESUME\_\allowbreak{}REJECTED} & Det/None \\
\texttt{V97} & Attach without a proof of possession is refused as possession\_\allowbreak{}invalid & \texttt{attach with possession none} & \texttt{(none: attach refused, possession\_\allowbreak{}invalid)} & no session \\
\texttt{V98} & Attach with a proof signed by another key is refused as possession\_\allowbreak{}invalid & \texttt{attach with possession wrong-key} & \texttt{(none: attach refused, possession\_\allowbreak{}invalid)} & no session \\
\texttt{V99} & Attach with a grant made for another public key under the same key id is refused as grant\_\allowbreak{}invalid & \texttt{attach with grant other-key} & \texttt{(none: attach refused, grant\_\allowbreak{}invalid)} & no session \\
\texttt{V100} & Attach with a grant made for another bridge is refused as grant\_\allowbreak{}invalid & \texttt{attach with grant other-bridge} & \texttt{(none: attach refused, grant\_\allowbreak{}invalid)} & no session \\
\texttt{V101} & Attach with a proof of possession made for another bridge is refused as possession\_\allowbreak{}invalid & \texttt{attach with possession other-bridge} & \texttt{(none: attach refused, possession\_\allowbreak{}invalid)} & no session \\
\texttt{V102} & Recovery repairs a registered public key that differs from the one in the log & \texttt{disconnect; tamper\_\allowbreak{}registry(ck-1,other-key); restart; reattach} & \texttt{ATTACHED, DETACHMENT, RESUME\_\allowbreak{}PREPARED, BRIDGE\_\allowbreak{}RECOVERED, DETACHMENT, RESUME\_\allowbreak{}PREPARED, RESUME\_\allowbreak{}STARTED, RESUME\_\allowbreak{}VERIFIED} & Att/None \\
\texttt{V103} & Recovery removes a registry entry the log does not record & \texttt{disconnect; tamper\_\allowbreak{}registry(ck-9,extra); restart} & \texttt{ATTACHED, DETACHMENT, RESUME\_\allowbreak{}PREPARED, BRIDGE\_\allowbreak{}RECOVERED, DETACHMENT, RESUME\_\allowbreak{}PREPARED} & Det/None \\
\texttt{V104} & Recovery accepts a registry that matches the log, with a key added by a connected client & \texttt{add\_\allowbreak{}client(ck-2); disconnect; restart; reattach(ck-2)} & \texttt{ATTACHED, CLIENT\_\allowbreak{}ADDED, DETACHMENT, RESUME\_\allowbreak{}PREPARED, BRIDGE\_\allowbreak{}RECOVERED, DETACHMENT, RESUME\_\allowbreak{}PREPARED, RESUME\_\allowbreak{}STARTED, RESUME\_\allowbreak{}VERIFIED} & Att/None \\
\texttt{V105} & A client registration without a signature is refused & \texttt{add\_\allowbreak{}client(ck-2,bad\_\allowbreak{}signature,none)!} & \texttt{ATTACHED} & Att/None \\
\texttt{V106} & A client registration signed by another holder of the key id is refused and the stranger cannot join & \texttt{add\_\allowbreak{}client(ck-2,bad\_\allowbreak{}signature,forged)!; join(ck-2)!} & \texttt{ATTACHED, RESUME\_\allowbreak{}REJECTED} & Att/None \\
\texttt{V107} & A signature made for a join is refused as a client registration & \texttt{add\_\allowbreak{}client(ck-2,bad\_\allowbreak{}signature,other-operation)!} & \texttt{ATTACHED} & Att/None \\
\texttt{V108} & A client registration presented a second time is refused & \texttt{add\_\allowbreak{}client(ck-2); add\_\allowbreak{}client(ck-2,nonce\_\allowbreak{}used,replay)!} & \texttt{ATTACHED, CLIENT\_\allowbreak{}ADDED} & Att/None \\
\texttt{V108b} & A client registration is verified before the connection is answered: an unregistered authorising key is a bad signature, a registered one that is not connected is via\_\allowbreak{}not\_\allowbreak{}connected & \texttt{add\_\allowbreak{}client(ck-2); add\_\allowbreak{}client(ck-3,bad\_\allowbreak{}signature)!; add\_\allowbreak{}client(ck-3,via\_\allowbreak{}not\_\allowbreak{}connected)!; add\_\allowbreak{}client(ck-3)} & \texttt{ATTACHED, CLIENT\_\allowbreak{}ADDED, CLIENT\_\allowbreak{}ADDED} & Att/None \\
\texttt{V109} & A challenge for a terminated session is refused exactly as one for an unknown session & \texttt{terminate; challenge\_\allowbreak{}uniform; challenge!} & \texttt{ATTACHED, TERMINATION} & Terminated \\
\texttt{V110} & A re-attach request for a terminated session is refused exactly as one for an unknown session & \texttt{terminate; reattach\_\allowbreak{}uniform(unknown); reattach(unknown)!} & \texttt{ATTACHED, TERMINATION} & Terminated \\
\texttt{V111} & Revoking by compromise the key whose re-attach is being verified abandons the re-attach & \texttt{disconnect; reattach(no-complete); revoke(ck-1,compromise); complete\_\allowbreak{}resume!} & \texttt{ATTACHED, DETACHMENT, RESUME\_\allowbreak{}PREPARED, RESUME\_\allowbreak{}STARTED, KEY\_\allowbreak{}REVOKED, DETACHMENT, RESUME\_\allowbreak{}PREPARED} & Det/None \\
\texttt{V111b} & Revoking by rotation the only key while its re-attach is verified abandons the re-attach; the session then ends at T\_\allowbreak{}detach & \texttt{disconnect; reattach(no-complete); revoke(ck-1,rotation); complete\_\allowbreak{}resume!; advance(600)} & \texttt{ATTACHED, DETACHMENT, RESUME\_\allowbreak{}PREPARED, RESUME\_\allowbreak{}STARTED, KEY\_\allowbreak{}REVOKED, DETACHMENT, RESUME\_\allowbreak{}PREPARED, DETACH\_\allowbreak{}TIMEOUT, TERMINATION} & Terminated \\
\texttt{V111c} & Revoking by rotation the key whose re-attach is verified abandons it; another key then resumes & \texttt{add\_\allowbreak{}client(ck-2); disconnect; reattach(ck-1,no-complete); revoke(ck-1,rotation); reattach(ck-2)} & \texttt{ATTACHED, CLIENT\_\allowbreak{}ADDED, DETACHMENT, RESUME\_\allowbreak{}PREPARED, RESUME\_\allowbreak{}STARTED, KEY\_\allowbreak{}REVOKED, DETACHMENT, RESUME\_\allowbreak{}PREPARED, RESUME\_\allowbreak{}STARTED, RESUME\_\allowbreak{}VERIFIED} & Att/None \\
\texttt{V112} & An approval of the content alone does not release a call under another tool name & \texttt{escalate; tool\_\allowbreak{}call(k1); answer(k1,content-only); resolve} & \texttt{ATTACHED, HUMAN\_\allowbreak{}PENDING, TOOL\_\allowbreak{}CALL\_\allowbreak{}QUEUED, TOOL\_\allowbreak{}CALL\_\allowbreak{}ANSWERED, HUMAN\_\allowbreak{}RESOLVED, TOOL\_\allowbreak{}CALL\_\allowbreak{}REJECTED} & Att/None \\
\texttt{V113} & The forwarded request is the hashed tool, permission and content, also when forwarded again after a restart & \texttt{escalate; tool\_\allowbreak{}call(k1); resolve; restart; reattach} & \texttt{ATTACHED, HUMAN\_\allowbreak{}PENDING, TOOL\_\allowbreak{}CALL\_\allowbreak{}QUEUED, TOOL\_\allowbreak{}CALL\_\allowbreak{}ANSWERED, HUMAN\_\allowbreak{}RESOLVED, TOOL\_\allowbreak{}CALL\_\allowbreak{}DRAINED, BRIDGE\_\allowbreak{}RECOVERED, DETACHMENT, RESUME\_\allowbreak{}PREPARED, RESUME\_\allowbreak{}STARTED, RESUME\_\allowbreak{}VERIFIED} & Att/None \\
\texttt{V114} & A session whose log cannot be replayed is quarantined alone; the others recover & \texttt{plant\_\allowbreak{}malformed\_\allowbreak{}log(bad-1,no-genesis); plant\_\allowbreak{}malformed\_\allowbreak{}log(bad-2,missing-field); restart; other\_\allowbreak{}quarantined(bad-1); other\_\allowbreak{}quarantined(bad-2)} & \texttt{ATTACHED, BRIDGE\_\allowbreak{}RECOVERED, DETACHMENT, RESUME\_\allowbreak{}PREPARED} & Det/None \\
\texttt{V114b} & An operator key addition is refused on a quarantined session & \texttt{revoke(ck-1,compromise); tamper\_\allowbreak{}log; restart; operator\_\allowbreak{}add\_\allowbreak{}key(ck-2)!; operator\_\allowbreak{}add\_\allowbreak{}key(ck-2,sig-none)!} & \texttt{ATTACHED, KEY\_\allowbreak{}REVOKED, DETACHMENT, RESUME\_\allowbreak{}PREPARED, RECOVERY\_\allowbreak{}FAILED} & Det/None \\
\texttt{V115} & A key-exhausted session keeps its detach timer and ends when T\_\allowbreak{}detach elapses & \texttt{revoke(ck-1,compromise); advance(599); advance(1)} & \texttt{ATTACHED, KEY\_\allowbreak{}REVOKED, DETACHMENT, RESUME\_\allowbreak{}PREPARED, DETACH\_\allowbreak{}TIMEOUT, TERMINATION} & Terminated \\
\texttt{V116} & A restart treats a key-exhausted session like any other: T\_\allowbreak{}detach restarts and the session ends & \texttt{revoke(ck-1,compromise); advance(400); restart; advance(599); advance(1)} & \texttt{ATTACHED, KEY\_\allowbreak{}REVOKED, DETACHMENT, RESUME\_\allowbreak{}PREPARED, BRIDGE\_\allowbreak{}RECOVERED, DETACHMENT, RESUME\_\allowbreak{}PREPARED, DETACH\_\allowbreak{}TIMEOUT, TERMINATION} & Terminated \\
\texttt{V117} & The operator adds a key before T\_\allowbreak{}detach elapses: the new key re-attaches and the timer stops & \texttt{revoke(ck-1,compromise); advance(300); operator\_\allowbreak{}add\_\allowbreak{}key(ck-2); reattach(ck-2); advance(5000)} & \texttt{ATTACHED, KEY\_\allowbreak{}REVOKED, DETACHMENT, RESUME\_\allowbreak{}PREPARED, OPERATOR\_\allowbreak{}ADD\_\allowbreak{}KEY, RESUME\_\allowbreak{}STARTED, RESUME\_\allowbreak{}VERIFIED} & Att/None \\
\texttt{V118} & After T\_\allowbreak{}detach a key-exhausted session is over: OPERATOR\_\allowbreak{}ADD\_\allowbreak{}KEY is refused & \texttt{revoke(ck-1,compromise); advance(600); operator\_\allowbreak{}add\_\allowbreak{}key(ck-2)!} & \texttt{ATTACHED, KEY\_\allowbreak{}REVOKED, DETACHMENT, RESUME\_\allowbreak{}PREPARED, DETACH\_\allowbreak{}TIMEOUT, TERMINATION} & Terminated \\
\texttt{V119} & A pending escalation on a key-exhausted session still times out & \texttt{escalate; revoke(ck-1,compromise); advance(599); advance(1)} & \texttt{ATTACHED, HUMAN\_\allowbreak{}PENDING, KEY\_\allowbreak{}REVOKED, DETACHMENT, RESUME\_\allowbreak{}PREPARED, HUMAN\_\allowbreak{}TIMEOUT, TERMINATION} & Terminated \\
\texttt{V120} & A key-exhausted session whose client stays connected has no timer until the transport is lost, then it ends after T\_\allowbreak{}detach & \texttt{revoke(ck-1,rotation); advance(5000); escalate!; disconnect; advance(599); advance(1)} & \texttt{ATTACHED, KEY\_\allowbreak{}REVOKED, DETACHMENT, RESUME\_\allowbreak{}PREPARED, DETACH\_\allowbreak{}TIMEOUT, TERMINATION} & Terminated \\
\texttt{V121} & A revocation signed for another bridge is refused and commits nothing & \texttt{revoke(ck-1,compromise,sig-other-bridge)!} & \texttt{ATTACHED} & Att/None \\
\texttt{V122} & A revocation signed at an earlier log position is refused and commits nothing & \texttt{add\_\allowbreak{}client(ck-2); revoke(ck-1,compromise,sig-stale)!} & \texttt{ATTACHED, CLIENT\_\allowbreak{}ADDED} & Att/None \\
\texttt{V123} & An operator key addition signed for another bridge is refused and the session stays key exhausted & \texttt{revoke(ck-1,compromise); operator\_\allowbreak{}add\_\allowbreak{}key(ck-2,sig-other-bridge)!} & \texttt{ATTACHED, KEY\_\allowbreak{}REVOKED, DETACHMENT, RESUME\_\allowbreak{}PREPARED} & Det/None \\
\texttt{V124} & An operator key addition signed at an earlier log position is refused and the session stays key exhausted & \texttt{revoke(ck-1,compromise); operator\_\allowbreak{}add\_\allowbreak{}key(ck-2,sig-stale)!} & \texttt{ATTACHED, KEY\_\allowbreak{}REVOKED, DETACHMENT, RESUME\_\allowbreak{}PREPARED} & Det/None \\
\texttt{V125} & An operator termination signed for another bridge is refused and the session stays key exhausted & \texttt{revoke(ck-1,compromise); operator\_\allowbreak{}terminate(sig-other-bridge)!} & \texttt{ATTACHED, KEY\_\allowbreak{}REVOKED, DETACHMENT, RESUME\_\allowbreak{}PREPARED} & Det/None \\
\texttt{V126} & An operator termination signed at an earlier log position is refused and the session stays key exhausted & \texttt{revoke(ck-1,compromise); operator\_\allowbreak{}terminate(sig-stale)!} & \texttt{ATTACHED, KEY\_\allowbreak{}REVOKED, DETACHMENT, RESUME\_\allowbreak{}PREPARED} & Det/None \\
\texttt{V127} & A re-attach on an attached session is refused as already\_\allowbreak{}attached and changes no cell & \texttt{add\_\allowbreak{}client(ck-2); reattach(ck-2)!} & \texttt{ATTACHED, CLIENT\_\allowbreak{}ADDED, RESUME\_\allowbreak{}REJECTED} & Att/None \\
\texttt{V128} & After a refused re-attach on an attached session, join is the way in & \texttt{add\_\allowbreak{}client(ck-2); reattach(ck-2)!; join(ck-2)} & \texttt{ATTACHED, CLIENT\_\allowbreak{}ADDED, RESUME\_\allowbreak{}REJECTED, CLIENT\_\allowbreak{}JOINED} & Att/None \\
\texttt{V129} & Recovery writes back a registry entry the log records & \texttt{add\_\allowbreak{}client(ck-2); disconnect; tamper\_\allowbreak{}registry(ck-2,missing); restart; reattach(ck-2)} & \texttt{ATTACHED, CLIENT\_\allowbreak{}ADDED, DETACHMENT, RESUME\_\allowbreak{}PREPARED, BRIDGE\_\allowbreak{}RECOVERED, DETACHMENT, RESUME\_\allowbreak{}PREPARED, RESUME\_\allowbreak{}STARTED, RESUME\_\allowbreak{}VERIFIED} & Att/None \\
\texttt{V130} & A bare terminate is refused on a key-exhausted session; only the operator's signed termination ends it & \texttt{revoke(ck-1); terminate!; operator\_\allowbreak{}terminate} & \texttt{ATTACHED, KEY\_\allowbreak{}REVOKED, OPERATOR\_\allowbreak{}TERMINATION, TERMINATION} & Terminated \\
\texttt{V131} & A re-attach naming an unregistered key is refused as a bad signature; the log keeps the true reason & \texttt{disconnect; reattach\_\allowbreak{}key\_\allowbreak{}uniform(ck-9)} & \texttt{ATTACHED, DETACHMENT, RESUME\_\allowbreak{}PREPARED, RESUME\_\allowbreak{}REJECTED, RESUME\_\allowbreak{}REJECTED} & Det/None \\
\texttt{V132} & Operator events refused on the wire with named reasons & \texttt{revoke(ck-9)!; operator\_\allowbreak{}add\_\allowbreak{}key(ck-2)!; operator\_\allowbreak{}terminate!; revoke(ck-1,compromise); operator\_\allowbreak{}add\_\allowbreak{}key(ck-1)!; revoke(ck-1,compromise)!} & \texttt{ATTACHED, KEY\_\allowbreak{}REVOKED, DETACHMENT, RESUME\_\allowbreak{}PREPARED} & Det/None \\
\texttt{V132b} & Operator events check the signature before they name the state of the session; a session that is over is not terminated again & \texttt{revoke(ck-9,sig-forged)!; operator\_\allowbreak{}add\_\allowbreak{}key(ck-2,sig-none)!; operator\_\allowbreak{}terminate(sig-forged)!; terminate; operator\_\allowbreak{}terminate(sig-forged)!; operator\_\allowbreak{}terminate!; revoke(ck-1)!} & \texttt{ATTACHED, TERMINATION} & Terminated \\
\texttt{V133} & A re-attach with a nonce issued for another session is refused and the owner then re-attaches & \texttt{attach\_\allowbreak{}other(ck-1,sess-133b); disconnect; reattach(ck-1,other-session)!; reattach(ck-1)} & \texttt{ATTACHED, DETACHMENT, RESUME\_\allowbreak{}PREPARED, RESUME\_\allowbreak{}REJECTED, RESUME\_\allowbreak{}STARTED, RESUME\_\allowbreak{}VERIFIED} & Att/None \\
\texttt{V134} & A client registration with a nonce issued for another session is refused & \texttt{attach\_\allowbreak{}other(ck-1,sess-134b); add\_\allowbreak{}client(ck-2,nonce\_\allowbreak{}unknown,other-session)!} & \texttt{ATTACHED} & Att/None \\
\bottomrule
\end{longtable}
\endgroup


\subsection{Event schemas}\label{sec:schemas}

Every event kind has a JSON Schema for its payload in
\texttt{supplementary/schemas/events/}, and \texttt{event.schema.json}
fixes the envelope (\texttt{kind}, \texttt{session\_id}, \texttt{sequence},
\texttt{prev\_hash}, \texttt{payload}, \texttt{timestamp}). The schemas close
the free-form gap that otherwise lets bridges satisfy the invariants and
still produce non-interoperable logs (Section~\ref{sec:limitations}). They
are generated by \texttt{schemas/build\_schemas.py}, and
\texttt{check\_schemas.py} validates every event the reference bridge emits
during the test suite against them and fails if an emitted kind has no
schema. \texttt{CHECKPOINT} is specified (I-CHECKPOINT, SHOULD) with
frozen field names (\texttt{root}, \texttt{tree\_size},
\texttt{signature}, \texttt{receipt}) and is emitted by the
reference bridge when the configured threshold is reached
(tests T63--T66, vectors V52--V54, V88, V89).

\subsection{Fault-injection and cost measurements}\label{sec:measurements}

The tests show that each invariant holds; they do not show what the
mechanism buys in operation or what it costs. The measurements of event loss and of verification cost address
that, produced by \texttt{supplementary/measure\_faults.py}
(\texttt{make measurements}), which is seeded and reports only counts
(calls, events, SHA-256 computations), so a rerun gives the same tables.

\paragraph{Event loss.} For each fault class the script generates
seeded schedules of tool calls, human escalations and answers, clock
advances and one injected fault type (client disconnects, human
timeouts, replays of re-attach requests the bridge had accepted, resubmitted call keys, bridge
restarts, and all of them together), and runs the same schedule against
the reference bridge and against an acknowledging relay. The relay is a
stand-in whose rules are listed in the script: it numbers every accepted
call and writes the call, its acknowledgement and the open prompt to a
store that survives a restart; it forwards while a client is attached and
no prompt is open, holds the other calls unacknowledged and replays them
on re-attach; it denies a prompt after $T_{\text{human}}$, drops the calls
behind it and records the drop; it never terminates a session on a
timeout; and it accepts any re-attach request, with no nonce and no
signature. It is the strongest baseline that still lacks an authenticated
resume, so the measurement asks what the contract adds beyond reliable
delivery. A call counts as \emph{accepted} if the agent was told so, as
\emph{lost} if it was neither forwarded nor left any recorded
disposition, and as \emph{orphaned} if it is in the log but the session
ended by timeout before it was forwarded, rejected or expired. Both
systems are scored on the same set: the calls both accepted and the replay
attempts both made, so the Accepted columns of a row are equal. (The
bridge refuses calls after a timeout has terminated the session; the relay
never terminates, so it accepts more calls than the bridge in those rows,
and the surplus is left out.)

% Generated by supplementary/measure_faults.py --emit-tex; do not edit.
\begin{table}[ht]\centering\small
\caption{Event loss under injected faults, an acknowledging relay against the bridge (seed 20260101, 200 schedules of 60 actions per row). Both systems are scored on the same calls and the same replay attempts. Accepted: calls both accepted. Lost: no delivery and no trace. Orphaned: queued and logged but without a recorded disposition (session terminated by a timeout). Dup: keys forwarded twice. Replays: re-attach requests that had been accepted, presented again, accepted / attempted. Early: requests that were refused before their key was registered, presented again before or after the registration, accepted / attempted.}\label{tab:fault-loss}
\begin{tabular}{@{}l l r r r r r r@{}}\toprule
Fault & System & Accepted & Lost (\%) & Orphaned & Dup & Replays & Early \\\midrule
no fault & relay & 6549 & 0 (0.0) & 0 & 0 & 0/0 & 0/0 \\
 & bridge & 6549 & 0 (0.0) & 0 & 0 & 0/0 & 0/0 \\
\addlinespace
disconnects & relay & 4393 & 0 (0.0) & 0 & 0 & 0/0 & 0/0 \\
 & bridge & 4393 & 0 (0.0) & 0 & 0 & 0/0 & 0/0 \\
\addlinespace
human timeouts & relay & 3565 & 0 (0.0) & 0 & 0 & 0/0 & 0/0 \\
 & bridge & 3565 & 0 (0.0) & 266 & 0 & 0/0 & 0/0 \\
\addlinespace
replays & relay & 3851 & 0 (0.0) & 0 & 906 & 837/837 & 0/0 \\
 & bridge & 3851 & 0 (0.0) & 0 & 714 & 0/837 & 0/0 \\
\addlinespace
pre-registration replays & relay & 3398 & 0 (0.0) & 0 & 0 & 0/0 & 1015/1015 \\
 & bridge & 3398 & 0 (0.0) & 0 & 0 & 0/0 & 0/1015 \\
\addlinespace
restarts & relay & 4854 & 0 (0.0) & 0 & 0 & 0/0 & 0/0 \\
 & bridge & 4854 & 0 (0.0) & 0 & 0 & 0/0 & 0/0 \\
\addlinespace
all faults & relay & 2133 & 0 (0.0) & 0 & 539 & 481/481 & 334/334 \\
 & bridge & 2133 & 0 (0.0) & 226 & 296 & 0/481 & 0/334 \\
\bottomrule
\end{tabular}\end{table}

\begin{table}[ht]\centering\small
\caption{Verification cost in SHA-256 computations against log length. Events: events of the plain log; the checkpointed log holds the same events plus the checkpoints and is the one verified. Full: recompute every event hash of it (I5). Suffix: recompute only from the latest checkpoint, worst and mean over all prefixes of the checkpointed log. Proof: one Merkle inclusion proof for a mid-log event (the leaf hash and one hash per level of the tree of RFC 9162). Writer: total cost of emitting the checkpoints (event hashes, leaf hashes and interior nodes), per event of the log; it grows with the log because the reference recomputes the whole tree at each checkpoint.}\label{tab:verify-cost}
\begin{tabular}{@{}r r r r r r r@{}}\toprule
Events & Checkpoints & Full & Suffix worst & Suffix mean & Proof & Writer per event \\\midrule
1003 & 1 & 1004 & 1000 & 499 & 11 & 3 \\
4003 & 4 & 4007 & 1002 & 501 & 13 & 7 \\
16003 & 16 & 16019 & 1002 & 502 & 15 & 25 \\
64003 & 64 & 64067 & 1002 & 502 & 17 & 97 \\
\bottomrule\end{tabular}\end{table}


Table~\ref{tab:fault-loss} shows no loss for either system on any
schedule: a relay that numbers and stores its calls, acknowledges them and
replays the unacknowledged ones loses nothing under disconnects, human
timeouts or restarts, so the bridge's log is not what separates the two
on this measure. The columns that differ are the others. Replays of
re-attach requests that the bridge had accepted are rejected in every
attempt by the bridge (nonce use and expiry, I9) and accepted in every
attempt by the relay, which has nothing to check them against; this is the
authentication the contract adds and the relay lacks by construction. The
measurement replays only requests that were accepted: a request refused
before its key was registered, which both bridges accept on replay once the
key is registered (limitations register), is not among its fault classes. The bridge forwards fewer duplicate
keys (the Dup column) because it holds a queued call under its key (I14);
the relay forwards a resubmitted key as a new call. The bridge does
\emph{not} remove duplicates in general: a call forwarded directly while a
client is attached and no prompt is open leaves no record, so a
resubmission under the same key is forwarded again. The cost of the
contract shows in the orphaned column: when a timeout terminates the
session, queued calls remain logged without a recorded disposition, which
an operator must treat as an audit finding; the relay avoids this by never
ending the session, which is the unbounded detachment the contract exists
to rule out.

\paragraph{Verification cost.} The script grows real logs of the
reference bridge to the lengths in Table~\ref{tab:verify-cost}, with and
without checkpoints, and counts SHA-256 computations. The checkpoints
are emitted by the count trigger at the default interval of $N=1000$
events since the previous checkpoint. Both logs are driven through the
same number of cycles, so they hold the same events; the checkpointed
log holds in addition the checkpoint events themselves (the
\emph{Events} column counts the events of the plain log, and the full
verification of the checkpointed log hashes those plus the
checkpoints). Full verification (I5) is linear in the log. Verifying
only the suffix after the latest checkpoint is bounded by the
checkpoint interval whatever the log length, and one Merkle inclusion
proof grows with $\log_2$ of the log. The price is on the writing side,
and it grows: the reference bridge recomputes the whole tree over all
preceding events at every checkpoint, so with a checkpoint every $N$
events the total writer cost is quadratic in the log length and the cost
per logged event grows linearly with it (last column: a few
computations per event at the shortest log, tens at the longest). This
growth belongs to the reference implementation, not to the contract:
I-CHECKPOINT requires the root over the preceding events and says
nothing about how it is computed, and an incremental (append-only) tree
that keeps only its right-hand frontier would make the work per
appended event constant and a checkpoint cost logarithmic in the log.
The reference does not have one, and the measurement does not test that
claim. The suffix figures assume the checkpoint root is
trusted; the signature field in the reference is a placeholder, and a verifier that recomputes the root
pays the full linear cost once.

\paragraph{Limits.} The same author wrote the bridge, the relay, the
fault injector and the loss definition, so a zero shows that the
mechanisms behave as designed on these schedules, not that they are
complete. The relay is this paper's own construction, chosen as the
strongest baseline without an authenticated resume; a different relay (one
that checks a nonce, or that does not record its drops) would give
different columns, and a real product was not measured. It records its
drops and survives a restart by assumption, which a deployment may not
provide. The faults are synthetic and independent of each other and the
schedules are short. Scoring on the calls both systems accepted makes
the columns comparable inside a row; rows with human timeouts and restarts
differ in the number of calls because the bridge terminates sessions, so the
rows are not comparable as rates. The loss measure says nothing about
approval: a relay that replays held calls after a disconnect may forward
a call no human answered, and \texttt{Lost} does not count that. The Rust
crate is not measured: it does not emit \texttt{CHECKPOINT}. Timing is not
reported because wall-clock figures of a single-process reference say little
about a deployment. The writer cost reported is that of the reference's
recompute-everything checkpoint, not a lower bound for the contract.

\paragraph{Conformance classes.} Following IETF conventions, the
conformance classes are:

\begin{itemize}[leftmargin=2em]
\item \textbf{MUST} (mandatory): Invariants $I_1, \ldots, I_{15}$,
I11-REVOKE, I-GRANT, I-OPERATOR and I-SIGN-JWKS; the corresponding tests and vectors. A bridge that fails any
MUST is non-conformant.
\item \textbf{SHOULD}: Invariants I-SIGN-EVENT (event signing in
GEC-less deployments, per GAR \S 6.1) and I-CHECKPOINT (anchored
checkpoints).
\item \textbf{MAY}: Custom event payload schemas; additional states or
conditions beyond the cells of the product state machine
(Definition~\ref{def:sm}). Extensions MUST be documented and MUST
NOT weaken any MUST invariant.
\end{itemize}

% ================================================================================
